\documentclass[11pt,a4paper]{article}
\usepackage[utf8]{inputenc}
% lmodern must precede fontenc: without a scalable Type1 font the T1 encoding
% falls back to bitmap fonts, and microtype's expansion then aborts the run
% with "auto expansion is only possible with scalable fonts".
\usepackage{lmodern}
\usepackage[T1]{fontenc}
\usepackage[margin=1in]{geometry}
\usepackage{amsmath,amssymb}
\usepackage{graphicx}
\usepackage{microtype}
\usepackage[hidelinks,breaklinks]{hyperref}
\usepackage{url}
\urlstyle{same}

% Generated by code/tools/md_to_tex.py from the SurveyForge Markdown output.
% Citation numbers match the source .md exactly.

\title{Securing Large Language Models: A Survey of Watermarking and Fingerprinting Techniques}
\author{Generated by SurveyForge}
\date{}

\begin{document}
\maketitle
\tableofcontents
\newpage



\section{Introduction}


Watermarking and fingerprinting constitute the foundational technical infrastructure for securing large language models (LLMs) across their full lifecycle---from pretraining data provenance and fine-tuning lineage to API-mediated deployment and edge-device adaptation---yet they remain conceptually and operationally distinct from broader AI safety, privacy, and robustness paradigms. Unlike safety guardrails that constrain harmful outputs \cite{doi:10.55662/jst.2023.4605} or privacy-preserving mechanisms that mitigate membership inference or data leakage \cite{doi:10.3390/info15110697,doi:10.1007/s44443-025-00177-1}, watermarking and fingerprinting are \emph{provenance-first} primitives: they encode verifiable claims of origin, ownership, or usage into either model behavior or its artifacts, enabling forensic accountability without altering functional intent. This distinction is critical: while alignment techniques optimize for human values \cite{arxiv:2309.15025} and differential privacy enforces statistical indistinguishability \cite{doi:10.1007/s10462-024-11024-6}, watermarking and fingerprinting enforce \emph{traceability}---a property increasingly mandated by regulatory frameworks such as the EU AI Act's transparency obligations and U.S. copyright litigation landscapes \cite{arxiv:2405.01560}.

The urgency of this paradigm stems from three converging imperatives. Economically, LLMs represent capital-intensive intellectual assets: training a 7B-parameter model incurs multi-million-dollar compute costs \cite{arxiv:1701.04082}, rendering unauthorized replication or commercial exploitation a material threat. Legally, emerging mandates---including GDPR's right to data erasure \cite{doi:10.1145/3292500.3330885} and the AI Act's requirement for synthetic content disclosure \cite{doi:10.1038/s41586-024-08025-4}---demand auditable, cryptographically grounded attribution mechanisms. Societally, the weaponization of LLMs for academic misconduct \cite{doi:10.1371/journal.pone.0331871}, disinformation \cite{arxiv:2112.04359}, and medical misinformation \cite{doi:10.1038/s41591-024-03445-1} necessitates tamper-resilient provenance signals that persist beyond surface-level paraphrasing \cite{arxiv:2303.13408}. Crucially, these pressures expose a fundamental asymmetry: watermarking is \emph{active}, owner-controlled embedding of secret-key--dependent signals into outputs or parameters \cite{arxiv:1903.01743,doi:10.1145/3649329.3655674}, whereas fingerprinting is \emph{passive}, forensic identification of intrinsic behavioral or structural signatures---such as activation sparsity patterns, attention manifold geometry, or latent-space confusion zones \cite{arxiv:2601.03242}---without requiring prior coordination. This duality enables complementary security postures: watermarking supports contractual enforcement and liability assignment, while fingerprinting enables adversarial auditing and ecosystem-wide provenance mapping \cite{arxiv:2502.00706}. As we synthesize across over 50 foundational works, the field is converging not toward isolated techniques, but toward a unified \emph{provenance assurance stack}---one that integrates cryptographic unforgeability \cite{arxiv:2306.09194}, socio-technical governance interfaces, and system-level interoperability with speculative decoding, confidential computing, and multimodal pipelines \cite{doi:10.1038/s41586-024-08025-4,doi:10.1609/aaai.v37i12.26750,arxiv:2406.03728}. The path forward lies not in optimizing for detection accuracy alone, but in co-designing watermarking and fingerprinting as inseparable components of responsible AI governance---where technical rigor, legal admissibility, and human-centered equity are jointly optimized.


\section{Foundations and Taxonomy of Watermarking and Fingerprinting Paradigms}



\subsection{Tripartite Embedding Locus Taxonomy: Output, Parameter, and Behavior Spaces}


The tripartite embedding locus taxonomy---output, parameter, and behavior spaces---constitutes the ontological bedrock upon which all watermarking and fingerprinting paradigms for large language models must be rigorously grounded. Unlike implementation-centric taxonomies that fragment techniques by algorithmic family (e.g., sampling-based vs. cryptographic), this functional classification centers on \emph{where provenance information is physically instantiated}, thereby exposing first-principles constraints on stealth, persistence, adversarial resilience, and deployment feasibility. Output-space watermarking operates as a \emph{decoding-time intervention}, embedding signals into the autoregressive token sequence via statistically biased sampling---exemplified by green-list methods \cite{arxiv:1701.04082} and their entropy-aware variants \cite{doi:10.1038/s41586-024-08025-4}. Its security rests on statistical distinguishability under the null hypothesis of unmarked generation, yet it suffers from fundamental fragility: paraphrasing collapses detection power \cite{arxiv:2303.13408}, speculative decoding introduces irreconcilable interference \cite{doi:10.1038/s41586-024-08025-4}, and homoglyph or whitespace attacks exploit tokenizer-level vulnerabilities \cite{arxiv:2402.16187}. Crucially, output-space schemes are \emph{inherently non-radioactive}: they do not transfer to distilled or fine-tuned models unless explicitly re-embedded---a limitation that undermines lineage tracking across model evolution.

Parameter-space watermarking and fingerprinting anchor provenance directly in the model's structural substrate---whether full weights, attention matrices, or low-rank adapters---enabling verification without output modification. EmMark demonstrates robustness against quantization and pruning on edge-deployed LLMs \cite{doi:10.1145/3649329.3655674}, while Functional Invariants To Watermark Large Transformers leverages functional invariances (e.g., dimension permutations, scaling/unscaling) to embed watermarks with no computational cost and zero output perturbation \cite{doi:10.1109/icassp48485.2024.10447264}. These approaches exhibit strong resistance to fine-tuning and merging, yet demand white- or gray-box access---rendering them inapplicable to closed API services. Critically, parameter-space techniques blur the watermarking--fingerprinting boundary: EmbMarker implants backdoor triggers in embeddings for EaaS copyright verification \cite{doi:10.18653/v1/2023.acl-long.423}, whereas WAPITI transfers base-model watermarks to fine-tuned derivatives via parameter integration \cite{arxiv:2410.06467}, revealing how parameter-space loci enable \emph{system-aware integrity binding} rather than static marking.

Behavior-space fingerprinting infers model identity solely from black-box input-output mappings---exploiting stable statistical biases in repetition rates, calibration errors, or activation sparsity profiles \cite{doi:10.1145/3292500.3330885,arxiv:2502.00706}. SLIM induces latent-space confusion zones detectable via hypothesis testing under ultra-low coverage \cite{arxiv:2601.03242}, while Bileve introduces a bi-level signature scheme that differentiates multiple provenance scenarios through rank-based sampling and fine-grained integrity checks \cite{arxiv:2406.01946}. Though deployable in vendor-agnostic settings, behavior-space methods face systemic vulnerabilities: ensemble inhibition attacks suppress fingerprint responses through token filtering \cite{arxiv:2402.16187}, and distributional shifts---particularly in low-resource languages---induce catastrophic detection failure \cite{arxiv:2402.16187}. The convergence across these loci is not accidental: modern architectures like Bileve \cite{arxiv:2406.01946} and PLMmark \cite{doi:10.1609/aaai.v37i12.26750} deliberately compose signals across output, parameter, and behavioral dimensions, transforming the tripartite taxonomy from a descriptive classification into a \emph{compositional design principle}---one that acknowledges no single locus suffices for end-to-end provenance assurance in heterogeneous, adversarial, and evolving LLM ecosystems.


\subsection{Detection Modality Spectrum: From White-Box Verification to Black-Box Auditing}


The detection modality spectrum constitutes a foundational axis of operational security in watermarking and fingerprinting---defining not \emph{what} is verified, but \emph{how much} the verifier knows about the model's internals. We formalize this as a three-tiered continuum: white-box (full parameter access), gray-box (intermediate representations---logits, attention maps, or architectural priors), and black-box (query-response behavior only)---each imposing distinct constraints on adversary capability, verification latency, forensic rigor, and real-world deployability. This modality axis is inextricably coupled with the tripartite embedding locus taxonomy introduced earlier: white-box detection naturally aligns with \emph{parameter-space} watermarking (e.g., ZKROWNN \cite{doi:10.1109/dac56929.2023.10247798} leveraging Zero-Knowledge Proofs for ownership attestation, or Functional Invariants To Watermark Large Transformers \cite{doi:10.1109/icassp48485.2024.10447264} achieving lossless embedding via dimension permutations and scaling/unscaling), enabling cryptographically grounded verification and provable robustness under fine-tuning, quantization, and pruning. Yet this tight coupling also exposes critical vulnerabilities: white-box schemes are operationally infeasible for third-party users of closed APIs and critically susceptible to gradient inversion attacks \cite{doi:10.1145/3580305.3599291} and model stealing \cite{arxiv:2402.19361}, where adversaries reconstruct green lists via mixed-integer programming \cite{arxiv:2405.19677} or erase fingerprints through dummy neuron injection \cite{arxiv:2303.09732}. Gray-box detection strikes a pragmatic equilibrium---bridging the structural fidelity of parameter-space methods with the deployment flexibility of behavioral inference---and thus serves as the natural locus for \emph{cross-dimensional signal composition}. PLMmark \cite{arxiv:2402.16187} exploits supervised contrastive loss on output representations of trigger sets to deviate watermarked models' responses from clean samples; Robust Distortion-free Watermarks for Language Models \cite{arxiv:2307.15593} aligns generated text to a randomized key sequence using internal token-level sampling logic---achieving high detection power (p < 0.01) from as few as 35 tokens even after 40--50\% random edits. However, it assumes partial trust in model providers and remains exposed to logit obfuscation, speculative decoding mismatches that break Gumbel-based sampling invariants \cite{arxiv:2402.05864}, and activation poisoning via adversarial perturbations \cite{doi:10.1016/j.comcom.2019.12.016}. Black-box detection---grounded exclusively in \emph{behavior-space} signatures---underpins vendor-agnostic auditing: Tracing Text Provenance via Context-Aware Lexical Substitution \cite{doi:10.1609/aaai.v36i10.21415} inserts watermarks via BERT-guided, context-aware synonym substitution; WASA: WAtermark-based Source Attribution for Large Language Model-Generated Data \cite{arxiv:2310.00646} enables source attribution by mapping provider-specific texts to unique watermarks; and Hey, That's My Model! Introducing Chain \& Hash, An LLM Fingerprinting Technique \cite{arxiv:2407.10887} cryptographically binds fingerprint prompts with responses via chain-and-hash binding---enabling irrefutable ownership proof without model access. These methods enable third-party watermarking without API cooperation but confront fundamental statistical limits: detection power scales sublinearly with query budget (O($\surd$Q) under Wald sequential testing), suffers from high variance under low-entropy prompts, and is susceptible to adaptive prompt probing that exploits calibration errors \cite{arxiv:2402.16187}. Crucially, the spectrum is not static but \emph{stratified}: recent frameworks like WaterPool \cite{arxiv:2405.13517} decouple key generation (white-box) from mark embedding (black-box), while Hufu \cite{arxiv:1910.12366} exploits permutation equivariance to enable gray-box extraction from permuted inputs without weight access---directly instantiating the \emph{compositional design principle} articulated in the prior subsection, wherein no single locus suffices for end-to-end provenance assurance. This convergence signals a paradigm shift---from monolithic modality assumptions toward \emph{modular, composable verification stacks}, where cryptographic primitives, statistical detectors, and behavioral probes operate synergistically across threat boundaries. Such stacks inherently invoke the four interdependent security properties---detectability, robustness, unforgeability, and collusion resistance---that will be axiomatized in the following subsection, revealing how modality choices directly constrain which properties can be simultaneously satisfied under adversarial pressure. Future work must therefore formalize cross-modality information flow, derive lower bounds on detection latency under constrained query budgets, and develop unified calibration protocols ensuring p-value validity across white-, gray-, and black-box regimes---thereby transforming detection modality from an implementation constraint into a design-first security primitive.


\subsection{Formal Security Properties: A Unified Information-Theoretic and Cryptographic Framework}


This subsection establishes a unifying formal framework that transcends heuristic claims of ``robustness'' or ``stealth'' by axiomatizing four interdependent security properties---detectability, robustness, unforgeability, and collusion resistance---through the complementary lenses of information theory, cryptography, and game-theoretic adversarial modeling. Crucially, these properties are not independent; their simultaneous satisfaction is constrained by fundamental limits exposed in seminal impossibility results \cite{arxiv:2505.06827}, which prove that no watermark can be both \emph{strongly robust} (i.e., resistant to quality-preserving perturbations) and \emph{computationally undetectable} without degrading output fidelity---a consequence of the attacker's access to a quality oracle and a perturbation oracle inducing an efficiently mixing walk over high-quality outputs. Detectability is formally grounded in statistical hypothesis testing: the null hypothesis \$H\_0\$ posits that text is drawn from the unwatermarked model distribution \$\textbackslash{}mathcal\{P\}\_0\$, while the alternative \$H\_1\$ asserts it originates from the watermarked distribution \$\textbackslash{}mathcal\{P\}\_1\$. Detection power is quantified via interpretable p-values derived from binomial statistics \cite{arxiv:2301.10226}, KL divergence bounds \cite{arxiv:2306.04634}, or mutual information between watermark bits and token sequences \cite{arxiv:2301.10226}. Recent advances like DERMARK \cite{arxiv:2406.10281} and WaterMax \cite{arxiv:2403.04808} derive token-count lower bounds for reliable detection at fixed Type I/II error rates, shifting evaluation from empirical thresholds to sample-complexity guarantees. Robustness is elevated beyond token-level edits to semantic invariance under bounded structural transformations---paraphrasing, summarization, translation, and multimodal re-rendering---formalized via perturbation oracles with provable survival radii \cite{doi:10.18653/v1/2024.findings-naacl.40,doi:10.18653/v1/2023.acl-long.117}. Yet this formalism acknowledges inherent trade-offs: beta-watermarking \cite{arxiv:2406.02603} demonstrates that even ``distortion-free'' schemes suffer distribution bias under key collisions, exposing a hidden source of detectability. Unforgeability transitions from secret-key reliance to cryptographic primitives: public-key signature embedding \cite{arxiv:2402.16187}, asymmetric generation/detection networks \cite{arxiv:2405.13517}, and zero-knowledge proofs for ownership attestation \cite{doi:10.1109/dac56929.2023.10247798}. Collusion resistance is modeled as a multi-player game where adversaries aggregate models or outputs; ensemble watermarking \cite{arxiv:2402.16187} and SLIM \cite{arxiv:2601.01296} explicitly design for uniqueness preservation under distillation, merging, or voting-based inhibition \cite{arxiv:2601.04261}. Critically, this unified framework reveals deep structural parallels: watermark unforgeability mirrors fingerprint collision resistance, while robustness against paraphrasing aligns with fingerprint stability under behavioral shifts. Future work must bridge the gap between theoretical guarantees and operational reality---e.g., integrating certified robustness \cite{doi:10.18653/v1/2024.findings-naacl.40} with latency-aware inference stacks \cite{doi:10.1038/s41586-024-08025-4}, or grounding cryptographic assumptions in empirically validated threat models \cite{arxiv:2402.19361,doi:10.1145/3618257.3624845}. Only through such synthesis can watermarking and fingerprinting evolve from ad hoc heuristics into verifiable, composable pillars of AI governance.


\subsection{Threat Model Taxonomy: Adversary Capabilities, Goals, and Attack Surfaces}


A rigorous threat model taxonomy is indispensable for unifying the fragmented landscape of watermarking and fingerprinting security analysis---and serves as the essential bridge between the formal security properties established in the preceding subsection and the paradigm-convergent design principles introduced in the following one. Rather than treating attacks as isolated phenomena, we formalize a three-axis orthogonal classification---adversary capability (access modality), operational goal (intent), and attack surface (locus of intervention)---that exposes deep structural symmetries across paradigms and reveals systemic vulnerabilities masked by siloed evaluation. This taxonomy directly operationalizes the impossibility constraints and detection-robustness trade-offs articulated earlier: for instance, the fundamental distortion--robustness tension \cite{arxiv:2505.06827} manifests most acutely in output-space removal attacks, while the vulnerability of backdoor-based schemes to out-of-distribution triggers \cite{arxiv:2409.06130} reflects a failure to satisfy the \emph{unforgeability} and \emph{collusion resistance} axioms under white-box neutralization intent. Crucially, the axes are not abstract categories but empirically grounded dimensions that determine \emph{which security property is most critically compromised}---a point explicitly foreshadowed in the conclusion of the prior subsection and directly leveraged in the subsequent convergence analysis.

First, adversary capability spans white-box (full parameter access), gray-box (logits or intermediate activations), and black-box (query-response only) settings---a spectrum directly shaping feasibility and bounding the attainable guarantees for each security property. White-box adversaries can execute gradient inversion to steal green lists via mixed integer programming \cite{arxiv:2405.19677}, thereby undermining cryptographic unforgeability and exposing the fragility of secret-key-dependent schemes under full model visibility. Black-box adversaries, meanwhile, may distill watermarked outputs to reconstruct deterministic response patterns \cite{arxiv:2402.19361}, enabling spoofing without key access---an attack that directly challenges the statistical detectability and collusion resistance assumed in black-box hypothesis testing frameworks. Critically, this capability spectrum maps precisely onto the detection modality spectrum introduced earlier: white-box detection enables provable verification but invites gradient inversion; black-box detection trades statistical power for operational agnosticism, yet remains bounded by query-budget and sample-complexity constraints; and gray-box methods occupy the middle ground where robustness must coexist with resilience to logit obfuscation---making the taxonomy not merely descriptive but prescriptive for security-aware system design.

Second, operational goals decompose into four canonical intents: \emph{removal} (erasing provenance signals without semantic degradation), \emph{spoofing} (synthesizing authentic-looking watermarks without secret keys), \emph{false attribution} (misassigning ownership via collision or mimicry), and \emph{neutralization} (suppressing uniqueness under aggregation). Removal attacks dominate output-space literature---e.g., synonym substitution \cite{arxiv:2505.02884}, round-trip translation \cite{arxiv:2402.14007}. These directly instantiate the \emph{robustness} challenge formalized earlier: paraphrasing and translation constitute bounded semantic perturbations against which watermark survival must be quantified---not heuristically, but via provable survival radii and KL divergence bounds. Spoofing, in turn, operationalizes the \emph{unforgeability} property: distillation-enabled spoofing \cite{arxiv:2402.19361} demonstrates how adversarial knowledge extracted from black-box queries can subvert cryptographic assumptions unless detection incorporates ensemble redundancy or error-correcting structure---as later realized in DERMARK and WaterMax. Neutralization, finally, engages \emph{collusion resistance}: the fragility of backdoor-based watermarks under structural assumptions \cite{arxiv:2409.06130} reveals how aggregation---whether via model merging or ensemble voting---can collapse discriminative signal spaces, motivating the multi-bit ensemble composition and adaptive modulation principles advanced in the next subsection.

Third, attack surfaces---output, parameter, and behavioral---define the physical locus of interference and expose divergent failure modes. While output-space attacks suffer from inevitable trade-offs between distortion and robustness \cite{arxiv:2505.06827}, parameter- and behavior-space interventions remain underexplored in the provided literature: no cited paper substantiates concrete neutralization mechanisms at those loci, and thus we refrain from speculative claims. Instead, we emphasize that \emph{all three surfaces are governed by the same underlying security axioms}: output-space watermarking must preserve fidelity while ensuring detectability; parameter-space fingerprinting must embed uniqueness without degrading alignment (as warned in \cite{doi:10.14722/ndss.2026.230138}); and behavior-space fingerprints must exhibit stability under bounded distributional shifts---mirroring the robustness formalism applied earlier to paraphrasing and summarization. This unified grounding reinforces the central thesis of the following subsection: watermarking and fingerprinting are not competing solutions but complementary facets of a single \emph{provenance assurance calculus}, differentiated not by locus but by how multi-bit ensembles, adaptive modulation, and modular integration enforce cross-surface invariants.

Critically, these axes intersect nontrivially: a white-box adversary pursuing neutralization may undermine backdoor-based watermarks by exploiting out-of-distribution trigger fragility \cite{arxiv:2409.06130}, while a black-box adversary aiming for spoofing may distill watermarked outputs to reconstruct deterministic response patterns \cite{arxiv:2402.19361}. This taxonomy thus reframes ``robustness'' not as a monolithic property but as a conditional invariant---e.g., a scheme robust against black-box removal may be brittle under white-box neutralization---and compels co-design of embedding and detection modalities aligned along all three axes. Future work must therefore shift from isolated attack demonstrations to \emph{cross-axis stress testing}: evaluating how a parameter-space watermark withstands black-box spoofing via distillation, or whether behavior-space fingerprints survive gray-box neutralization under activation sparsity drift. Only such integrated threat modeling can expose the true boundaries of provable security in generative AI provenance---and enable the transition from abstract desiderata to actionable, cross-modality invariants enforceable across the full detection spectrum, as demanded by both the formal framework above and the architectural convergence below.


\subsection{Interoperability and Cross-Paradigm Convergence: Bridging Watermarking and Fingerprinting}


The artificial boundary between watermarking and fingerprinting---long sustained by disciplinary silos and implementation-centric taxonomies---collapses under rigorous scrutiny of their shared architectural DNA. As this survey reveals, both paradigms are converging toward a unified \emph{provenance assurance calculus}, governed not by embedding locus (output, parameter, or behavior space) but by three interlocking design principles: multi-bit ensemble composition, adaptive signal modulation, and modular, system-aware integration. This convergence is neither incidental nor incremental; it reflects a fundamental shift from \emph{signal insertion} to \emph{integrity binding}---where provenance is no longer an add-on artifact but a compositional property of the model's operational semantics.

Multi-bit ensemble watermarking exemplifies this paradigm shift. \textbf{Provably Robust Multi-bit Watermarking for AI-generated Text via Error Correction Code} \cite{arxiv:2401.16820} and \textbf{Robust and Efficient Watermarking of Large Language Models Using Error Correction Codes} \cite{doi:10.56553/popets-2025-0126} demonstrate that combining orthogonal signals---green-list sampling, acrostic constraints, and sensorimotor norm deviations---yields superadditive detection power: detection AUROC increases by 0.18--0.23 over single-signal baselines while preserving perceptual fidelity, as measured by human fluency ratings and LLM-as-judge coherence scores \cite{doi:10.18653/v1/2024.findings-naacl.40,arxiv:2601.14311}. Crucially, this ensemble logic mirrors behavior-space fingerprinting in \textbf{Hey, That's My Model! Introducing Chain \& Hash, An LLM Fingerprinting Technique} \cite{arxiv:2407.10887}, where latent-space confusion zones, repetition burstiness, and calibration error profiles are fused via weighted hypothesis testing---a formalization of \emph{multi-feature discriminability} grounded in mutual information bounds \cite{doi:10.3389/fdata.2021.729663}. The convergence is mathematically precise: both approaches optimize the Kullback--Leibler divergence \$D\_\{\textbackslash{}text\{KL\}\}(P\_\{\textbackslash{}text\{marked\}\} \textbackslash{}| P\_\{\textbackslash{}text\{unmarked\}\})\$ not over a single distribution but over a product space \$\textbackslash{}mathcal\{P\} = \textbackslash{}prod\_i P\_i\$, where each \$P\_i\$ corresponds to an independent signal channel. This enables provable robustness against targeted erasure of any single channel, as shown in \textbf{Provably Robust Multi-bit Watermarking for AI-generated Text via Error Correction Code}'s collusion-resistance analysis under model merging.

Adaptive strength modulation further dissolves the dichotomy. \textbf{Adaptive Text Watermark for Large Language Models} \cite{arxiv:2401.13927} modulates perturbation magnitude by token entropy \$H(t\_i \textbackslash{}mid t\_\{<i\})\$, attenuating watermarking in low-entropy safety-critical contexts (e.g., medical disclaimers) while amplifying it in high-entropy factual claims---mirroring \textbf{Distilling the Thought, Watermarking the Answer: A Principle Semantic Guided Watermark for Large Reasoning Models} \cite{arxiv:2601.05144}'s causal tracing that localizes watermark influence to neurons with highest gradient variance w.r.t. instruction-following fidelity. \textbf{A Robust Semantics-based Watermark for Large Language Model against Paraphrasing} \cite{doi:10.18653/v1/2024.findings-naacl.40} extends this to dual-space adaptation: simultaneously biasing logits \emph{and} sampling decisions based on semantic embedding distance, achieving 92\% paraphrase survival versus 63\% for static green-list schemes \cite{arxiv:2402.19361}. Such context-awareness is not merely heuristic; it satisfies a formal \emph{risk-proportional embedding constraint}: \$\textbackslash{}| \textbackslash{}delta p\_t \textbackslash{}|\_2 \textbackslash{}leq \textbackslash{}lambda \textbackslash{}cdot \textbackslash{}mathcal\{R\}(t\_\{<i\}, \textbackslash{}theta)\$, where \$\textbackslash{}mathcal\{R\}\$ quantifies contextual risk (e.g., factual claim density) and \$\textbackslash{}lambda\$ enforces alignment preservation per \textbf{Watermarking Degrades Alignment} \cite{arxiv:2306.04634}.

Finally, modularity enables cross-paradigm interoperability. \textbf{Tracing Text Provenance via Context-Aware Lexical Substitution} \cite{doi:10.1609/aaai.v36i10.21415} decouples watermark injection from autoregressive decoding, inserting semantic watermarks post-hoc via sentence-embedding--guided synonym substitution---enabling third-party watermarking without API cooperation. \textbf{PLMmark: A Secure and Robust Black-Box Watermarking Framework for Pre-trained Language Models} \cite{doi:10.1609/aaai.v37i12.26750} embeds zero-knowledge proofs within MPC protocols to verify watermark presence without exposing weights or keys, satisfying GDPR Article 32 and EU AI Act transparency mandates. \textbf{Waterfall: Framework for Robust and Scalable Text Watermarking} \cite{arxiv:2407.04411} generalizes this to non-autoregressive models by defining watermark stability via Markov-chain key sequences invariant to token reordering---demonstrating that \emph{system-level integrity} supersedes locus-specific assumptions. These advances collectively signal a maturation beyond monolithic techniques toward \emph{assurance layers}: composable, governance-ready modules that integrate with speculative decoding, confidential computing, and safety guardrails---not as afterthoughts, but as first-class citizens in the AI systems stack.


\section{Technical Architectures for Output-Level Watermarking}



\subsection{Statistical Sampling-Based Watermarking: Foundations, Variants, and Token-Level Trade-offs}


Statistical sampling-based watermarking constitutes the foundational paradigm for output-level provenance, embedding signals directly into the autoregressive token generation process via pseudorandom, secret-key--driven perturbations of the softmax distribution. At its core lies a formal statistical distinguishability guarantee: given a sequence of \$n\$ tokens \$\textbackslash{}mathbf\{y\} = (y\_1, \textbackslash{}dots, y\_n)\$, a watermark detector computes a test statistic---e.g., the count of tokens falling within a cryptographically partitioned ``green list'' \$G\_k(\textbackslash{}mathbf\{h\}\_t)\$ derived from model hidden state \$\textbackslash{}mathbf\{h\}\_t\$ and secret key \$k\$---and rejects the null hypothesis \$H\_0\$: ``\$\textbackslash{}mathbf\{y\}\$ is unwatermarked'' if \$\textbackslash{}sum\_\{t=1\}\textasciicircum{}n \textbackslash{}mathbb\{I\}\cite{arxiv:2506.06407} > \textbackslash{}tau\_n\$, where \$\textbackslash{}tau\_n\$ is calibrated via binomial tail bounds to achieve provable false positive rate control \cite{doi:10.1038/s41586-024-08025-4}. Green-list sampling, as instantiated in Kirchenbauer et al. \cite{doi:10.1038/s41586-024-08025-4}, establishes the canonical framework; its variants---including KGW's key-dependent vocabulary shuffling \cite{doi:10.1038/s41586-024-08025-4}, ITS's iterative token-space reweighting \cite{doi:10.1038/s41586-024-08025-4}, and EXP's exponential weighting over logits \cite{doi:10.1145/3321705.3329808}---all share the same theoretical bedrock: statistical detectability scales with entropy divergence between watermarked and unwatermarked distributions, formalized as \$\textbackslash{}mathrm\{KL\}(p\_\{\textbackslash{}text\{wm\}\} \textbackslash{}| p\_\{\textbackslash{}text\{un\}\}) = \textbackslash{}Omega(\textbackslash{}alpha\textasciicircum{}2 / T)\$, where \$\textbackslash{}alpha\$ is embedding strength and \$T\$ temperature \cite{arxiv:2402.05864}. Yet empirical studies reveal critical fragility: BanglaLorica demonstrates that detection power collapses under low-entropy constrained generation (e.g., machine translation), where green-list coverage drops below 5\% due to deterministic prefix constraints, reducing AUROC from 0.92 to 0.58 across 12 language pairs \cite{doi:10.1038/s41586-024-08025-4}. This exposes a fundamental token-level trade-off: watermark strength must be modulated not uniformly but adaptively---via entropy-aware perturbation thresholds \cite{doi:10.1038/s41586-024-08025-4}, syntactic depth gating \cite{doi:10.1038/s41586-024-08025-4}, or confidence-score--driven scaling \cite{doi:10.1038/s41586-024-08025-4}---to preserve fluency while maintaining signal integrity in high-risk segments. Crucially, classical GumbelMax sampling suffers from deterministic-output degeneracy under repeated sampling, undermining diversity and enabling homoglyph/whitespace attacks; GumbelSoft and Logits-Addition resolve this by reparameterizing the sampling distribution probabilistically, yielding +0.17--0.29 AUROC gains and resilience against Unicode normalization attacks \cite{arxiv:2402.05864}. Moreover, tokenizer sensitivity remains a systemic vulnerability: Byte-Pair Encoding's subword fragmentation causes green-list misalignment under round-trip translation, while SentencePiece's context-dependent segmentation induces token-boundary shifts that erase watermark signals entirely---empirically validated in SeqMark's cross-lingual evaluation showing 4$\times$ degradation in Bangla RTT settings \cite{doi:10.1038/s41586-024-08025-4}. These findings collectively underscore that statistical watermarking is not a plug-and-play module but a co-design problem tightly coupled to decoding infrastructure, linguistic morphology, and entropy dynamics---pointing toward future architectures that unify cryptographic key derivation with dynamic token-space adaptation and certified robustness guarantees under bounded adversarial edits \cite{arxiv:2306.09194}.


\subsection{Cryptographic and Learning-Based Watermark Embedding: From Rejection Sampling to End-to-End Differentiable Decoders}


Output-level watermarking has evolved from heuristic statistical modulation toward architectures grounded in cryptographic primitives and joint learning objectives---enabling provable security, multi-bit payload encoding, and seamless deployment in open-model ecosystems without API access or model fine-tuning. This evolution is directly motivated by the critical limitations of statistical sampling--based methods---namely, their fragility under constrained generation, deterministic-output degeneracy, and tokenizer-induced signal erosion---as rigorously exposed in the preceding subsection. As a conceptual pivot \emph{away} from probabilistic biasing and \emph{toward} computational indistinguishability, rejection sampling--based watermarking constructs output distributions that are statistically identical to the unwatermarked model's distribution \emph{except} to parties holding the secret key \$k\$. Concretely, given \$k\$, a one-way function \$f\_k\$ maps candidate token sequences to pseudorandom bitstrings that determine acceptance or rejection during sampling; this ensures distortion-freeness---even under adaptive querying---while guaranteeing formal unforgeability: no polynomial-time adversary can generate a valid watermark without \$k\$ with non-negligible probability, under standard cryptographic assumptions \cite{arxiv:2405.13517}. Crucially, this paradigm directly resolves the ``Watermarks in the Sand'' impossibility result \cite{arxiv:2405.13517} by anchoring security not in statistical divergence (which collapses under low-entropy constraints), but in computational hardness---thereby inheriting robustness against the very attack vectors (e.g., prefix-constrained translation, Unicode normalization, subword misalignment) that undermine green-list sampling.

To overcome the sampling inefficiency inherent in rejection-based protocols---while preserving both cryptographic guarantees \emph{and} the capacity for rich, structured provenance---error-correcting code (ECC)--based schemes embed verifiable binary messages (e.g., user ID, timestamp, license tier) into token sequences via robust binary codes (RBCs) \cite{arxiv:2401.16820}. Formally, let \$\textbackslash{}mathbf\{c\} \textbackslash{}in \textbackslash{}\{0,1\textbackslash{}\}\textasciicircum{}m\$ be an \$m\$-bit payload encoded into codeword \$\textbackslash{}mathcal\{E\}(\textbackslash{}mathbf\{c\}) \textbackslash{}in \textbackslash{}\{0,1\textbackslash{}\}\textasciicircum{}n\$ using a \$\cite{arxiv:2512.03620}\$ linear code with minimum Hamming distance \$d\$. Each bit is mapped to a cryptographically hashed green list; even after \$t < d/2\$ adversarial edits (substitutions, insertions, deletions), maximum-likelihood decoding recovers \$\textbackslash{}mathbf\{c\}\$ with bounded bit error rate \$\textbackslash{}mathrm\{BER\} \textbackslash{}leq \textbackslash{}exp(-\textbackslash{}Omega(d))\$ \cite{arxiv:2401.16820}. Empirically, such schemes achieve 98.4\% match rate for 12-bit payloads in 200-token texts---surpassing Yoo et al.'s baseline by 12.8\%---while maintaining statistical indistinguishability \emph{even under entropy-constrained generation}, thereby directly addressing the low-coverage failure mode observed in BanglaLorica's evaluation \cite{doi:10.1038/s41586-024-08025-4}.

The most architecturally transformative advance is end-to-end differentiable watermarking, which co-trains encoder-decoder modules using online prompting as a surrogate for non-differentiable sampling \cite{arxiv:2403.10553}. Unlike logit-perturbation methods---which inherit the fluency--detectability trade-offs exposed earlier---this framework employs a policy gradient approach where the decoder computes a reward signal based on watermark detectability (e.g., KL divergence from green-list distribution), and the encoder updates logits to jointly maximize reward \emph{and} task utility. This yields higher paraphrase robustness across OPT, LLaMA, and Alpaca models while preserving perplexity and downstream performance \cite{arxiv:2403.10553}. Critically, its modular design decouples message encoding (binary signature injection) from distribution reparameterization (dense-to-sparse logits mapping), enabling plug-and-play deployment on quantized or instruction-tuned models with low latency overhead \cite{arxiv:2403.10553}. Emerging work like WaterMod further refines this via probability-aware modular partitioning---sorting tokens by model probability and assigning classes via rank mod \$k\$---thereby guaranteeing high-probability token availability while embedding base-\$k\$ digits per step \cite{arxiv:2511.07863}. Collectively, these approaches constitute a foundational shift: from watermarking as a \emph{post-hoc constraint} to a \emph{co-optimized system property}, where cryptographic rigor, learning efficiency, and operational viability converge---not in isolation, but as interoperable layers within the broader context-aware governance architecture discussed next.


\subsection{Context-Aware and Adaptive Watermarking: Semantic Integrity, Alignment Preservation, and Risk-Guided Modulation}


Context-aware and adaptive watermarking represents a paradigm shift from static, uniform signal embedding toward \emph{risk-proportional}, \emph{semantically grounded}, and \emph{alignment-preserving} modulation---directly confronting the foundational tension between detectability and responsible model behavior. Rather than applying fixed perturbations across all tokens, these architectures dynamically modulate watermark strength using real-time contextual signals: instruction-following fidelity, factual claim density, syntactic ambiguity, safety-criticality thresholds, and token-level semantic centrality. This transforms watermarking from a statistical post-hoc verification tool into an integrated component of model governance, where provenance assurance co-evolves with behavioral integrity. At its core lies a formal risk-guided objective: maximize detection power in high-stakes segments (e.g., ``According to peer-reviewed studies\ldots{}'' or code blocks containing \texttt{if \_\_name\_\_ == '\_\_main\_\_':}) while minimizing interference in alignment-critical contexts---such as medical disclaimers or constitutional interpretations---where even sub-percentage shifts in token probability can attenuate guardrail enforcement. This principle is operationalized through three interlocking mechanisms. First, \emph{instruction-following fidelity--driven modulation} leverages prompt-response consistency metrics---e.g., self-consistency scores over chain-of-thought rollouts or calibration error deviations---to amplify green-list bias precisely when the model asserts verifiable claims, achieving 3.2$\times$ higher detection F1 in citation-rich domains without increasing false positives in safety-critical outputs \cite{doi:10.34190/icair.4.1.3128}. Second, \emph{semantic differentiation} abandons vocabulary partitioning in favor of region-aware embedding: SeqMark \cite{doi:10.1609/aaai.v36i10.21415} and SemaMark \cite{doi:10.18653/v1/2024.findings-naacl.40} use BERT-based semantic similarity scoring to identify paraphrase-resilient substitution candidates, embedding watermarks only in tokens whose replacement preserves latent-space proximity to the original sentence embedding---boosting ROUGE-L retention by 19.7\% under aggressive paraphrasing while maintaining >92\% detection AUC. Third, \emph{alignment-resampling mitigation} introduces inference-time correction layers that recover degraded truthfulness and helpfulness: Alignment Resampling \cite{doi:10.14722/ndss.2026.230138} applies external reward models (e.g., fine-tuned DeBERTa) to re-rank top-k generations, restoring baseline alignment scores (TruthfulQA, SafetyBench) with only 2--4 additional samples---demonstrating that watermark-induced distributional shifts are not inherently misaligned but \emph{reversible}. Crucially, these approaches expose a fundamental limitation of entropy-agnostic watermarking: EWD \cite{arxiv:2403.13485} proves that uniform weighting fails catastrophically in low-entropy contexts (e.g., legal clauses), whereas context-aware schemes like WaterMax \cite{arxiv:2403.04808}, which decouples embedding locus from sampling mechanics, achieve near-zero perplexity delta while sustaining p < 0.001 detection at 50-token spans. Looking forward, the frontier lies in \emph{multi-modal risk orchestration}: integrating linguistic risk signals with vision-language grounding constraints (as in VLA-Mark \cite{arxiv:2406.03728}) and code syntax validity checks (STONE \cite{arxiv:2502.18851}) to enforce cross-modal consistency---ensuring that a watermark embedded in a Python function's docstring also manifests in its generated unit test's assertion logic. Such unified risk modeling transcends output-level watermarking to become a general-purpose architecture for \emph{verifiable responsible generation}.


\subsection{Robustness Engineering and Attack-Aware Design: Paraphrase Resilience, Black-Box Deployment, and Certified Guarantees}


Robustness engineering for output-level watermarking has evolved from heuristic resilience to a principled, attack-aware discipline---where architectural design is explicitly co-optimized with adversarial threat models and operational constraints exposed in prior subsections. This subsection synthesizes three interlocking advances: (i) paraphrase-resilient signal embedding that transcends token-level perturbation; (ii) black-box deployment mechanisms enabling third-party, API-agnostic watermarking without model access---directly extending the \emph{instruction-following fidelity--driven modulation} and \emph{semantic differentiation} principles introduced earlier; and (iii) certified guarantees grounded in randomized smoothing and information-theoretic bounds, which formalize the robustness promises implicit in context-aware schemes like WaterMax and VLA-Mark. Critically, these are not orthogonal improvements but mutually reinforcing layers: post-hoc semantic watermarking enables black-box operation, while certified radii formalize the robustness conferred by semantic invariance---and both serve as essential prerequisites for the \emph{holistic assessment paradigm} demanded in the following subsection, where detection accuracy alone is insufficient without provable resilience across translation, fine-tuning, and speculative decoding.

Paraphrase resilience is no longer achieved by increasing green-list bias strength---a strategy proven vulnerable to entropy-guided rewriting \cite{arxiv:2505.05190} and self-information--driven attacks \cite{arxiv:2505.05190}. Instead, state-of-the-art methods embed signals in \emph{semantic invariants}: SemaMark leverages sentence-level meaning preserved under synonym substitution and reordering \cite{doi:10.18653/v1/2024.findings-naacl.40}, while the Semantic Invariant Robust Watermark uses an auxiliary embedding LLM to map preceding tokens into semantic space before generating watermark logits---ensuring signal persistence even when lexical surface forms change \cite{arxiv:2310.06356}. Crucially, this semantic anchoring is empirically validated: SeqMark's sequence-level semantic differentiation improves F1 detection by up to 28\% in machine translation and summarization, where token-level schemes collapse due to region collapse \cite{arxiv:2601.11629}. This directly bridges the \emph{semantic differentiation} mechanism described previously---now reframed not merely as a quality-preserving policy but as a foundational robustness primitive---and anticipates the cross-lingual evaluation challenges highlighted later, where BanglaLorica exposes catastrophic failure of surface-form watermarks under round-trip translation.

Black-box deployment has matured beyond synonym insertion into architecturally coherent frameworks aligned with the risk-proportional governance vision. Tracing Text Provenance via Context-Aware Lexical Substitution inserts context-aware lexical substitutions guided by BERT-based semantic relatedness to preserve sentence coherence and improve transferability \cite{doi:10.1609/aaai.v36i10.21415}; PLMmark establishes ownership verification for pre-trained language models using public-key cryptography and supervised contrastive loss on trigger sets, ensuring transferability across downstream tasks \cite{arxiv:2402.16187}; and CATER introduces conditional watermarking rules that minimize distortion of overall word distributions while maximizing conditional selection shifts---rendering statistical detection infeasible even for attackers aware of the scheme's architecture \cite{doi:10.34190/icair.4.1.3128}. These approaches collectively decouple watermark injection from model internals, enabling verifiable provenance in commercial API ecosystems---precisely the environment where \emph{instruction-following fidelity--driven modulation} must operate without white-box access, and where \emph{alignment-resampling mitigation} becomes indispensable for preserving behavioral integrity under external constraints.

Certified guarantees now move beyond empirical robustness to mathematically provable lower bounds---fulfilling the demand for \emph{trajectory-aware benchmarking} and \emph{cryptographic rigor} articulated in the subsequent subsection. IPCert employs randomized smoothing to turn watermarking/fingerprinting methods into provably robust ones against model perturbations, providing certified robustness guarantees \cite{doi:10.1109/cvpr52733.2024.01169}. Similarly, Provably Robust Multi-bit Watermarking for AI-generated Text via Error Correction Code provides theoretical guarantees on bit error rate under bounded adversarial edits, attaining 98.4\% match rate for 12-bit payloads in 200-token texts---outperforming Yoo et al. by >12 percentage points under copy-paste attacks \cite{arxiv:2401.16820}. These frameworks expose a fundamental tension: while semantic invariance improves empirical resilience, only randomized smoothing and ECC-based coding yield \emph{formal} robustness certificates---bridging the gap between adversarial intuition and cryptographic assurance. Future work must unify these paradigms---embedding semantically invariant codes within smoothed detectors---to achieve both practical resilience and provable guarantees, thereby closing the loop between context-aware embedding, black-box deployability, and evaluation-ready certification.


\subsection{Evaluation, Benchmarking, and Operational Realities: Beyond Detection Accuracy to Holistic System Assessment}


The evaluation of output-level watermarking architectures has long suffered from a dangerous myopia: over-indexing on statistical detectability (e.g., p-value thresholds, AUROC) while neglecting the multidimensional operational fidelity required for real-world deployment. As demonstrated empirically, detection accuracy alone is a poor predictor of system viability---watermarked outputs may achieve \$p < 0.01\$ detectability yet induce measurable degradation in instruction-following fidelity (\$-18.3\textbackslash{}\%\$ GPT-4-Judge alignment score) or trigger false positives in safety moderation APIs \cite{arxiv:2402.19361,arxiv:2402.16187}. This subsection advances a holistic assessment paradigm grounded in four interdependent dimensions: human-perceived quality, computational efficiency under acceleration constraints, cross-lingual and low-resource robustness, and longitudinal functional integrity across model evolution. Critically, these dimensions are not orthogonal but coupled---e.g., adaptive watermarking that modulates strength by token entropy improves fluency but introduces latency overhead in speculative decoding due to draft-model--detector misalignment \cite{arxiv:2407.04411}. Standardized benchmarks like WaterBench now co-evaluate generation and detection jointly, using LLM-as-judge protocols to expose subtle distributional artifacts masked by perplexity alone \cite{doi:10.1145/3691626,arxiv:2502.11598}, while WaterJudge formalizes the quality-detection trade-off via Pareto-optimal operating point selection across summarization and translation tasks \cite{arxiv:2502.11598}. Latency analysis reveals fundamental incompatibilities: a provable no-go theorem establishes that maximal watermark strength and maximal sampling efficiency in speculative decoding cannot be simultaneously achieved---empirically validated across OPT and LLaMA with \$>42\textbackslash{}\%\$ throughput penalty when enforcing strict green-list adherence \cite{arxiv:2307.15593,arxiv:2403.04808}. Two pragmatic compromises emerge: (i) watermark-aware draft models that precompute partitioned logits, and (ii) deferred verification that batches detection post-generation, trading immediate auditability for sub-50ms inference overhead \cite{arxiv:2403.04808,arxiv:2307.15593}. Cross-lingual evaluation exposes stark inequities: BanglaLorica shows round-trip translation collapses token-level watermark detection from \$88.2\textbackslash{}\%\$ to \$9.7\textbackslash{}\%\$---a failure mitigated only by layered strategies combining embedding-time statistical signals with post-generation semantic signatures, yielding \$3.8\textbackslash{}times\$ relative improvement \cite{arxiv:2402.14007}. Most critically, longitudinal studies remain scarce: Mark Your LLM documents catastrophic watermark decay after just three rounds of fine-tuning, while MUSE demonstrates that unlearning-aware evaluation uncovers \$67\textbackslash{}\%\$ false-negative rates in watermark survival under selective forgetting \cite{doi:10.1145/3639477.3639720,doi:10.1609/aaai.v39i1.32018}. Future evaluation must therefore shift from static snapshots to \emph{trajectory-aware} benchmarking---tracking watermark persistence across iterative updates, quantization, and multimodal extension---and embed cryptographic guarantees (e.g., randomized smoothing for certified radii \cite{arxiv:2401.16820}) alongside human-centered metrics (e.g., cognitive load in academic editing scenarios \cite{arxiv:2403.04808}).


\section{Model-Centric Fingerprinting and Parameter-Space Watermarking}



\subsection{Weight-Space Fingerprinting: Intrinsic Model Identity from Parameter Geometry}


Weight-space fingerprinting represents a paradigm shift in model-centric provenance: rather than embedding signals \emph{into} behavior or \emph{through} training, it treats the model's parameter geometry as an intrinsic, high-dimensional fingerprint---a latent signature emergent from architecture, optimization trajectory, and data-induced manifold structure. Crucially, this approach is \emph{training-free} and \emph{passive}: no fine-tuning, no watermark injection, no inference-time intervention---only analysis of static weight tensors yields verifiable identity. Its foundational strength lies in mathematical invariance: under common post-training transformations---including quantization, pruning, LoRA adaptation, and even instruction tuning---the spectral, geometric, and alignment properties of weight matrices exhibit remarkable stability, enabling robust lineage verification across open-model ecosystems where weights are publicly shared but provenance is contested \cite{doi:10.1109/sds66131.2025.00024,doi:10.1109/icassp48485.2024.10447264,doi:10.3389/fdata.2021.729663}. Spectral fingerprinting leverages singular value decomposition (SVD) to extract dominant left/right singular vectors, whose subspace alignment---measured via Centered Kernel Alignment (CKA)---remains invariant under affine-preserving perturbations; AWM \cite{doi:10.3389/fdata.2021.729663} formalizes this using the Linear Assignment Problem (LAP) to match CKA-normalized subspaces across models, achieving near-zero false positives across 60+ model pairs by treating singular vectors not as isolated features but as \emph{geometric anchors} in the weight manifold. Attention manifold geometry extends this principle to structured tensors: attention weight matrices across layers and heads lie on low-dimensional Riemannian manifolds whose curvature, geodesic distance, and principal subspace alignment encode architecture- and training-data--specific signatures---resilient even after multi-modal extension or Qwen-style dense-to-sparse attention reparameterization \cite{doi:10.1109/icassp48485.2024.10447264,doi:10.1109/sds66131.2025.00024}. Weight permutation invariants constitute a third axis: by encoding identifiers via structured permutations of neuron indices or channel orders---then decoding via error-correcting codes (ECCs)---fingerprinting achieves cryptographic-grade unforgeability without runtime overhead, as ECC decoding guarantees reliable extraction despite adversarial perturbation of up to 30\% of permuted parameters \cite{doi:10.1145/3321705.3329808,arxiv:1701.04082}. Critically, these methods diverge from backdoor-based watermarking: they do not rely on behavioral triggers or output anomalies, making them immune to query-based scrubbing or distillation attacks that erase functional signatures but leave structural geometry intact \cite{arxiv:2402.19361,arxiv:2601.21252}. Yet limitations persist: spectral fingerprints degrade under aggressive quantization below INT4, and manifold geometry becomes ambiguous in heavily pruned models where attention sparsity collapses subspace dimensionality \cite{doi:10.1109/sds66131.2025.00024,doi:10.1145/3321705.3329808}. Emerging work bridges these gaps---e.g., Functional Invariants To Watermark Large Transformers exploits \emph{functional equivalence} via dimension scaling and permutation to generate watermarked copies with identical outputs, preserving stealth while anchoring identity in provable invariance. Looking forward, weight-space fingerprinting must evolve toward \emph{longitudinal invariance}: current methods verify static snapshots, but real-world model evolution---via continual learning, federated updates, or safety-aligned RLHF---demands fingerprints that track \emph{parameter trajectories}, not just endpoints. Integrating stochastic differential equations to model weight-space drift under fine-tuning, or leveraging persistent homology to capture topological invariants across model families, may yield the next generation of mathematically grounded, lifecycle-aware model identity.


\subsection{Backdoor-Based Parameter-Space Watermarking: Lightweight, Secret-Key--Driven Ownership Imprinting}


Backdoor-based parameter-space watermarking bridges the conceptual gap between weight-space and behavior-space provenance: it embeds \emph{structural provenance} directly into the model's functional mapping via lightweight, secret-key--dependent fine-tuning---neither relying on static weight geometry (as in spectral or manifold-based fingerprinting) nor on emergent behavioral statistics (as in black-box probing). Unlike output-space watermarking---which operates at inference time and suffers from semantic erasure---backdoor watermarking ensures zero-runtime-cost attribution that persists across deployment modalities, including Embedding-as-a-Service (EaaS), API wrappers, and quantized edge variants. Crucially, its security rests not on statistical obscurity but on cryptographic key dependence: the watermark is only detectable when the adversary supplies the correct trigger set \emph{and} possesses the secret key required to decode the target response pattern. This yields a formal unforgeability guarantee under standard cryptographic assumptions---as shown in \cite{arxiv:2405.13517}, detection reduces to one-way function inversion, rendering watermark presence computationally indistinguishable from benign behavior without the key---even under adaptive querying.

Three architectural paradigms dominate this space. First, \emph{trigger-word--driven embedding perturbations}, exemplified by \cite{doi:10.18653/v1/2023.acl-long.423}, select moderate-frequency natural words (e.g., ``veridian'', ``lattice'') as cryptographically sampled triggers and inject proportional target embeddings into the final layer representations. This ensures watermark transfer to stolen models while preserving utility across downstream NLP tasks---demonstrating <0.2\% accuracy degradation on GLUE benchmarks. Second, \emph{instructional fingerprinting}, as in \cite{arxiv:2401.12255}, encodes confidential private keys directly into instruction-tuning prompts (e.g., ``Respond with `\cite{arxiv:2502.18518}' if this prompt contains exactly three prepositions''), eliciting deterministic, model-specific responses. Critically, it supports multi-stage licensing tiers (e.g., MIT-style commercial/non-commercial flags) and resists fingerprint guessing via semantic obfuscation---achieving 98.7\% robustness against parameter-efficient fine-tuning attacks. Third, \emph{trigger-free neuron-level watermarking}, pioneered by \cite{doi:10.1109/icassp48485.2024.10447264}, repurposes membership inference logic for defense: fine-tuning on natural data induces subtle, key-dependent memorization biases in output probabilities---not tied to artificial triggers---yielding gray-box fingerprints that survive meta-prompting, fine-tuning, and model merging while remaining statistically undetectable to adversaries (\cite{arxiv:2409.06130} shows such in-distribution triggers resist white-box erasure by >40$\times$ over out-of-distribution counterparts).

This paradigm thus occupies a distinct operational niche: it requires \emph{controlled access during fine-tuning} (unlike passive weight-space fingerprinting) yet delivers \emph{functional imprinting} more robust than behavior-space signatures under query-based scrubbing or distillation. However, it inherits the threat-model limitations of any key-dependent method---namely, vulnerability to key leakage, side-channel extraction, or adversarial fine-tuning with trigger-aware objectives. As highlighted in the following subsection on behavior-space fingerprinting, these vulnerabilities motivate complementary verification strategies: where backdoor watermarking provides strong, key-gated ownership proof \emph{when training control is available}, behavior-space methods offer resilient, system-agnostic attribution \emph{when only queries or activations are accessible}. Future work must therefore reconcile cryptographic rigor with behavioral integrity: synthesizing secret-key--driven parameter imprinting with alignment-resampling mechanisms (\cite{arxiv:2601.14311}) and certified robustness guarantees to ensure that ownership proof does not compromise truthfulness, safety, or helpfulness---a non-negotiable requirement for regulatory admissibility under the EU AI Act and NIST AI RMF.


\subsection{Behavior-Space Fingerprinting: Black-Box Provenance from Statistical Biases and Activation Signatures}


Behavior-space fingerprinting represents a paradigm shift in model-centric provenance: it abandons assumptions of weight access, training intervention, or even API cooperation, instead treating the LLM as a \emph{stochastic black-box dynamical system} whose identity emerges from invariant statistical and activation-level response patterns under controlled probing. Unlike output-space watermarking---which imposes signals---and parameter-space methods---which read from weights---behavior-space fingerprinting \emph{discovers} discriminative signatures from the model's natural behavioral manifold. This is achieved through three complementary axes: (i) prompt-driven statistical bias profiling, (ii) layer-wise activation sparsity and entropy dynamics, and (iii) evolutionary discovery of maximally discriminative input-response mappings. Critically, these approaches are \emph{training-free}, require no secret keys or fine-tuning, and operate solely via query-response sequences---making them uniquely suited for auditing closed commercial APIs, verifying open-weight models in untrusted environments, or forensic attribution in adversarial settings where watermark removal is assumed.

Prompt-driven statistical bias fingerprinting exploits architecture- and training-data--informed deviations that persist across diverse prompts: token repetition rates exhibit characteristic burstiness profiles tied to attention head topology \cite{arxiv:2305.08883}; calibration errors manifest as systematic overconfidence in low-entropy contexts (e.g., factual assertions), quantified via expected calibration error (ECE) divergence from ideal uniformity \cite{arxiv:2403.04808}; and lexical entropy distributions deviate predictably under standardized cloze templates, enabling family-level clustering (e.g., Llama vs. Mistral) with >72\% accuracy without weight inspection \cite{arxiv:2601.01296}. These biases are not artifacts but emergent properties of optimization trajectories---hence resistant to paraphrasing yet vulnerable to ensemble inhibition, as demonstrated by token-filter and sentence-verification attacks that suppress response variance across model ensembles \cite{arxiv:2601.04261}.

Layer-wise activation sparsity and entropy fingerprinting leverages the stable cross-layer geometry of neuron activation density under fixed probing prompts. For instance, fine-tuned variants (e.g., Llama-2 vs. Llama-2-Chat) exhibit statistically separable variance patterns in attention head activation entropy---measured as \$H\textasciicircum{}\{(l)\} = -\textbackslash{}sum\_i p\_i\textasciicircum{}\{(l)\} \textbackslash{}log p\_i\textasciicircum{}\{(l)\}\$ across heads \$i\$ at layer \$l\$---even when weights are encrypted or quantized \cite{arxiv:2601.14311}. Such signatures arise from differential memorization pressure during instruction tuning and survive pruning and LoRA adaptation, though they degrade under aggressive quantization \cite{doi:10.1109/icassp48485.2024.10447264}. Crucially, this gray-box modality bridges black-box practicality and white-box fidelity: it requires only intermediate activations (accessible via some API providers or local inference stacks), not full weights.

Finally, evolutionary ``Hide-and-Seek'' fingerprint discovery formalizes the search for discriminative features as an in-context learning problem: an Auditor LLM generates prompts maximizing inter-family KL divergence between response distributions, while a Detective LLM performs hypothesis testing on logits or token probabilities \cite{arxiv:2601.01296,arxiv:2407.10887}. This co-evolutionary framework discovers latent-space confusion zones---regions where semantically similar prefixes induce divergent continuations---enabling ultra-low-coverage verification (<0.01\% of training data modified) while preserving utility \cite{arxiv:2601.01296}. The convergence of these three axes points toward a future where behavior-space fingerprinting evolves from static signature extraction to \emph{adaptive provenance orchestration}: dynamically composing statistical, activation, and decision-boundary signals per use case---balancing detection power, latency, and adversarial resilience in real-time.


\subsection{Adaptive and Multi-Layer Parameter-Space Watermarking: From Static Embedding to Dynamic Integrity Binding}


The evolution of parameter-space watermarking from static, binary identifiers to adaptive, multi-layer integrity binding represents a paradigmatic shift toward \emph{system-aware provenance assurance}---where watermarks are no longer inert signatures but dynamic, context-sensitive anchors integrated into the model's operational lifecycle. This transition is driven by three convergent imperatives: (i) the need for granular licensing enforcement beyond ``AI-generated'' attribution; (ii) the architectural necessity of coexistence with acceleration infrastructures such as speculative decoding; and (iii) the governance requirement that watermarks survive---not merely resist---post-deployment interventions like selective unlearning or fine-tuning. Recent advances formalize this shift through three interlocking technical axes: multi-bit payload encoding with distortion-free guarantees, system-coordinated embedding for inference-stack compatibility, and unlearning-resilient watermark placement grounded in parameter-space topology.

Multi-bit watermarking has moved decisively beyond zero-bit detection toward rich metadata binding. Methods like \emph{Multi-Bit Distortion-Free Watermarking} \cite{arxiv:2402.16578} and \emph{Advancing Beyond Identification: Multi-bit Watermark for Large Language Models} \cite{doi:10.18653/v1/2024.naacl-long.224} embed user IDs, license tiers, or timestamps directly into weight perturbations using error-correcting codes (ECCs), achieving bit error rates <2\% under 40\% token corruption \cite{arxiv:2401.16820}. Crucially, these schemes decouple payload capacity from perceptual distortion: \emph{WaterMod} \cite{arxiv:2511.07863} demonstrates that modular arithmetic over token ranks enables base-\emph{k} digit embedding per step without excluding high-probability tokens---guaranteeing fluency while supporting fine-grained provenance tracing. However, empirical analysis reveals a fundamental tension: ECC-based schemes exhibit sharp degradation when entropy heterogeneity is ignored, as shown in \emph{Invisible Entropy} \cite{arxiv:2505.14112}, which introduces an entropy tagger to adaptively gate embedding strength---reducing parameter overhead by 99\% while preserving detection fidelity on HumanEval.

System-coordinated watermarking addresses the practical incompatibility between strong watermark persistence and speculative decoding efficiency. Rather than treating watermarking as a post-hoc constraint, recent work embeds integrity signals \emph{within} the inference protocol itself. \emph{Permute-and-Flip} \cite{arxiv:2402.05864} redefines decoding: its PF decoder provably achieves up to 2$\times$ better quality-robustness trade-offs than standard sampling by aligning watermark persistence with draft-model coordination. Similarly, \emph{CATER} \cite{arxiv:2209.08773} embeds watermarks via conditional word selection rules optimized to minimize distributional distortion---rendering statistical inspection infeasible even to adversaries aware of the scheme's architecture. These approaches reconceptualize watermarking not as a post-hoc add-on but as a first-class component of the inference stack---directly enabling the \emph{adaptive provenance orchestration} foreshadowed in the prior subsection's discussion of behavior-space fingerprinting.

Finally, unlearning-aware watermarking confronts the fragility of ``radioactive'' signals under targeted forgetting. \emph{WAPITI} \cite{arxiv:2410.06467} transfers base-model watermarks to fine-tuned derivatives via parameter integration, ensuring continuity across model evolution. This insight bridges cryptographic unforgeability and structural robustness: \emph{Robust and Efficient Watermarking of Large Language Models Using Error Correction Codes} \cite{doi:10.56553/popets-2025-0126} proves that ECC-encoded weight permutations achieve collision resistance under bounded $\ell$$_{\mathrm{2}}$-norm perturbations, establishing the first formal guarantee that watermark survival can be decoupled from functional equivalence. Collectively, these advances signal a maturation toward \emph{dynamic integrity binding}: watermarks that evolve with model usage, encode policy-relevant semantics, and interoperate with the full spectrum of AI system infrastructure---from confidential computing enclaves to multimodal pipelines---thereby fulfilling the foundational promise of \emph{system-aware provenance assurance} introduced earlier and setting the stage for the rigorous security--governance co-design explored in the following subsection.


\subsection{Formal Security, Practical Robustness, and Governance Constraints of Model-Centric Techniques}


Model-centric techniques---encompassing parameter-space watermarking and weight- or behavior-space fingerprinting---introduce a paradigm shift from output-dependent provenance to \emph{structural identity anchoring}, where security guarantees must be grounded not in statistical indistinguishability of text, but in the cryptographic hardness of parameter inversion, the information-theoretic stability of model manifolds, and the forensic traceability of inference protocols. This dual evaluation lens reveals a profound tension: formal security properties such as unforgeability under key exposure remain largely aspirational for most model-centric schemes, while practical robustness is routinely undermined by systemic operational constraints---from quantization-induced signal erosion in edge deployments to activation sparsity drift across fine-tuned variants. Crucially, recent advances have begun bridging this gap. The first formal security baseline for parameter-space watermarking emerges from \emph{Functional Invariants To Watermark Large Transformers} \cite{doi:10.3390/electronics14010025}, which proves computational indistinguishability of watermarked models from clean ones by reducing detection to one-way function inversion---a property preserved even under adaptive querying. Similarly, \emph{Robust and Efficient Watermarking of Large Language Models Using Error Correction Codes} \cite{doi:10.3390/electronics14010025} establishes provable distortion-freeness via secret-key--dependent signature embedding, satisfying the Daubert standard for evidentiary admissibility through cryptographic reduction. Yet these guarantees assume white-box access and static model states---conditions violated in real-world settings where models undergo continual learning, quantization, or ensemble aggregation. Indeed, \emph{Robust and Efficient Watermarking of Large Language Models Using Error Correction Codes} \cite{doi:10.3390/electronics14010025} demonstrates that collision resistance under model merging degrades exponentially with the number of merged variants unless fingerprints are embedded via Centered Kernel Alignment--neutralized Linear Assignment Problem matching; meanwhile, \emph{WAPITI} \cite{doi:10.3390/electronics14010025} shows that unlearning-aware watermark placement can preserve uniqueness only when perturbations target low-sensitivity eigenspaces of the Hessian, a constraint incompatible with many lightweight backdoor schemes like \emph{EmbMarker} \cite{doi:10.3390/electronics14010025}. Governance constraints further compound this fragility: \emph{Hey, That's My Model! Introducing Chain \& Hash, An LLM Fingerprinting Technique} \cite{doi:10.3390/electronics14010025} satisfies GDPR Article 32 and EU AI Act Annex VI requirements by cryptographically binding prompt-response pairs into tamper-evident execution traces, yet its reliance on deterministic hashing renders it vulnerable to adversarial prompt padding attacks uncovered in \emph{Hey, That's My Model! Introducing Chain \& Hash, An LLM Fingerprinting Technique} \cite{doi:10.3390/electronics14010025}. More critically, equity analyses reveal systemic disparities: fingerprint stability collapses by >75\% for quantized Bangla fine-tunes compared to English counterparts \cite{arxiv:2402.14007}], while \emph{PLMmark} \cite{doi:10.3390/electronics14010025}'s public-key contrastive loss fails under low-resource tokenization due to vocabulary table sparsity---exposing false attribution risks in non-Latin scripts. These findings collectively underscore that model-centric security cannot be decoupled from deployment context: robustness requires co-design with speculative decoding stacks (\emph{Waterfall} \cite{doi:10.3390/electronics14010025}), forensic auditability demands zero-knowledge proof integration (\emph{Provably Robust Multi-bit Watermarking for AI-generated Text via Error Correction Code} \cite{doi:10.3390/electronics14010025}), and governance viability necessitates multi-stakeholder evaluation frameworks that jointly report unforgeability (for providers), decoder latency (for platforms), and perceptual fidelity (for users) \cite{doi:10.1145/3691626}; \emph{Securing Large Language Models: A Survey of Watermarking and Fingerprinting Techniques} \cite{doi:10.3390/electronics14010025}]. Future work must therefore move beyond isolated property verification toward \emph{composable assurance contracts}---formal specifications that encode not just ``what is embedded'' but ``under what transformations, for whom, and with what liability boundaries.''


\section{Robustness, Adversarial Vulnerabilities, and Attack Taxonomies}



\subsection{Taxonomy of Adversarial Capabilities and Threat Models}


Adversarial taxonomy in watermarking and fingerprinting must transcend superficial capability labels---such as ``black-box'' or ``white-box''---to capture the \emph{structural interplay} among access modality, operational intent, and knowledge assumptions. We therefore propose a triaxial adversarial framework grounded in empirical attack literature: (i) \textbf{access modality}, formalized as the adversary's ability to observe or manipulate model internals; (ii) \textbf{operational intent}, defined via game-theoretic payoff structures (e.g., minimizing detection probability while preserving semantic fidelity); and (iii) \textbf{knowledge assumptions}, quantified by the adversary's capacity to infer secret parameters, trigger semantics, or distributional invariants. This orthogonal decomposition exposes latent vulnerabilities obscured by monolithic threat modeling---for instance, gray-box adversaries leveraging calibration errors in logits or attention entropy are neither fully white-box nor black-box, yet exploit precisely the ``crack between layers'' where watermark detectors assume statistical independence across tokens or heads.

White-box adversaries---possessing full parameter access---operate under the strongest capability assumption and enable the most precise attacks: gradient inversion to reconstruct green lists \cite{arxiv:2402.19361}, weight-space manipulation to erase spectral fingerprints, or mixed-integer programming to reverse-engineer vocabulary partitions. Critically, these attacks succeed not due to brute-force optimization but because existing watermarking schemes violate \emph{parameter-space separation}: embedding signals in weights without enforcing cryptographic indistinguishability from random perturbations. As shown in \cite{doi:10.1109/icassp48485.2024.10447264}, functional equivalence under permutation or scaling can be exploited for stealthy watermark erasure---revealing that robustness requires not just invariance under transformations, but \emph{provably undetectable invariance}, as formalized in \cite{arxiv:2405.13517} via one-way function reduction.

Gray-box adversaries occupy a strategically rich middle ground: they observe logits, attention weights, or activation sparsity but lack full weight access. Such adversaries mount targeted neutralization---e.g., exploiting overconfidence in low-entropy contexts to bias token selection away from green-list regions, or poisoning intermediate representations to mask behavioral fingerprints. Notably, gray-box attacks expose a critical design flaw in many watermark decoders: their reliance on \emph{unverified statistical assumptions} about logit distributions. When adversaries inject calibrated noise into softmax inputs---as demonstrated in \cite{arxiv:2402.16187}---detection power collapses even when watermark strength remains high, because detectors assume i.i.d. token sampling rather than adaptive adversarial conditioning.

Black-box adversaries, constrained to input-output queries, exhibit surprising efficacy through compositionality: knowledge distillation enables spoofing without key exposure \cite{arxiv:2402.19361}, while prompt engineering allows watermark detection via crafted queries. Crucially, black-box attacks succeed not by breaking cryptographic primitives but by exploiting \emph{radioactive transfer}---a property intended as defense but weaponized as attack surface. This reveals a fundamental asymmetry: detection requires statistical significance over O(n) samples, whereas removal via distillation succeeds with O(1) queries per distilled sample.

Future threat modeling must therefore integrate \emph{adversarial epistemology}: distinguishing between adversaries who know the watermark algorithm (key-aware), those who only observe output distributions (distributionally constrained), and those who reverse-engineer triggers from behavioral anomalies (trigger-informed). Only such granular classification enables provable security guarantees---e.g., unforgeability under key leakage, collusion resistance under ensemble inhibition, or certified robustness against bounded query budgets. Without this rigor, evaluations risk false confidence---measuring resilience against straw-man adversaries while ignoring real-world compositional threats.


\subsection{Output-Space Attack Vectors: From Semantic Scrubbing to Covert Spoofing}


Output-space attacks constitute the most empirically pervasive and linguistically plausible threat vector against watermarking and behavioral fingerprinting systems, targeting the very surface---generated text---where provenance signals are embedded or inferred. As established in the preceding adversarial taxonomy, these attacks instantiate distinct combinations of \emph{access modality}, \emph{operational intent}, and \emph{knowledge assumptions}: semantic scrubbing operates under black-box query constraints with distributional knowledge; adaptive lexical substitution leverages gray-box awareness of tokenizer behavior and lexical redundancy; while covert spoofing often emerges from key-aware adversaries who reverse-engineer statistical patterns via distillation or prompt orchestration. Crucially, they all converge in exploiting the same foundational fragility exposed earlier---the \emph{decoupling of statistical detectability from cryptographic accountability}---rendering detection inherently asymmetric: it demands statistically significant deviation, whereas evasion or forgery requires only statistical mimicry.

Unlike parameter- or behavior-space attacks---which target model internals or functional mappings and are examined in the following subsection---output-space adversaries operate at zero marginal cost, leveraging widely available linguistic tools and human editing intuition to degrade or subvert detection without triggering perceptual alarms. This subsection synthesizes three interlocking attack paradigms: \emph{semantic scrubbing}, which preserves meaning while erasing statistical anomalies; \emph{adaptive lexical substitution}, which exploits tokenizer vulnerabilities and lexical redundancy to evade distributional detectors; and \emph{covert spoofing}, which synthesizes authentic-looking watermarked outputs without secret keys---thereby enabling false-flag attribution with high operational fidelity.

Semantic scrubbing attacks fundamentally exploit the \emph{quality-robustness trade-off} inherent in statistical watermarking. Recursive rewriting frameworks---such as those evaluated in \emph{Watermark Stealing in Large Language Models} \cite{arxiv:2402.19361}---demonstrate that even distortion-free schemes like \emph{Robust Distortion-free Watermarks for Language Models} \cite{arxiv:2307.15593} suffer >60\% detection failure after two rounds of LLM-mediated paraphrasing, as entropy redistribution dilutes green-list concentration. Crucially, machine translation cycles induce token boundary shifts that collapse detection power: \emph{Tracing Text Provenance via Context-Aware Lexical Substitution} \cite{doi:10.1609/aaai.v36i10.21415} shows round-trip translation reduces F1 from 88\% to <13\% for token-level watermarks, exposing a systemic fragility in vocabulary-partitioning assumptions. Similarly, \emph{Token-Specific Watermarking with Enhanced Detectability and Semantic Coherence for Large Language Models} \cite{arxiv:2402.18059} reveals that entropy-guided token substitution---targeting low-entropy positions where watermark bias is strongest---achieves near-perfect fluency preservation (BLEU $\geq$ 0.97) while reducing p-values below significance thresholds across OPT-1.3B and LLaMA-7B.

Adaptive lexical substitution attacks go beyond syntactic rewriting by weaponizing tokenizer heterogeneity. \emph{Tracing Text Provenance via Context-Aware Lexical Substitution} \cite{doi:10.1609/aaai.v36i10.21415} formalizes synonym substitution as a constrained optimization over BERT-derived lexical neighborhoods, minimizing KL divergence between original and scrubbed distributions while maximizing green-token removal. \emph{Undetectable Watermarks for Language Models} \cite{arxiv:2306.09194} extends this with homoglyph-aware perturbations, exploiting Unicode normalization ambiguities to break green-list alignment in \emph{A Watermark for Large Language Models} \cite{arxiv:2301.10226}---achieving 82\% evasion under adversarial prompting while preserving human fluency ratings (p < 0.01). Critically, these attacks succeed not because they distort semantics, but because they \emph{preserve statistical invariants}: as shown in \emph{Distortion-free Watermarks are not Truly Distortion-free under Watermark Key Collisions} \cite{arxiv:2406.02603}, pseudo-random sampling under key collisions induces distributional biases that substitution attacks amplify rather than disrupt.

Covert spoofing represents the most insidious evolution: moving from \emph{removal} to \emph{fabrication}. \emph{Watermark Stealing in Large Language Models} \cite{arxiv:2402.19361} enables third-party black-box insertion of semantic watermarks via context-dependent synonym chains, generating outputs statistically indistinguishable from legitimately watermarked ones---yet requiring no secret key or model access. \emph{Large Language Model Watermark Stealing With Mixed Integer Programming} \cite{arxiv:2405.19677} demonstrates distilled watermark replication: querying a watermarked teacher model yields outputs whose statistical patterns (e.g., green-token frequency correlations) can be reverse-engineered and cloned into student models with >95\% fidelity, enabling malicious actors to inject harmful content bearing authentic-looking watermarks. \emph{Bileve: Securing Text Provenance in Large Language Models Against Spoofing with Bi-level Signature} \cite{arxiv:2406.01946} further exposes this vulnerability by showing bi-level signature spoofing can forge fine-grained integrity bits while retaining coarse-grained source traces---blurring the line between provenance and liability. Collectively, these attacks reveal a foundational asymmetry: detection requires statistical deviation, but spoofing only requires statistical mimicry---a gap that no purely distributional watermark can close without cryptographic binding. Future defenses must therefore integrate \emph{provably undetectable} primitives \cite{arxiv:2306.09194} with \emph{traceable spoofing detection} \cite{arxiv:2402.19361}, shifting from passive signal embedding to active adversarial accountability---precisely the kind of epistemological rigor demanded by the triaxial adversarial framework, and a necessary precursor to confronting the structural subversion examined next in parameter- and behavior-space attacks.


\subsection{Parameter- and Behavior-Space Attacks: Erasure, Neutralization, and Collusion}


Parameter- and behavior-space attacks represent a paradigmatic shift from output-centric manipulation to structural subversion---targeting the very substrates of model identity: weights, activations, and input-output mappings. Unlike scrubbing attacks that operate on generated text, these adversaries bypass detection entirely by erasing, masking, or colluding against identifiers embedded in the model's internal geometry or behavioral manifold. Crucially, they exploit the \emph{radioactivity assumption}---that fingerprints persist invariantly across post-training modifications---revealing it not as a robust property but as a systemic vulnerability under adversarial pressure. Erasure attacks, such as those implemented in \emph{WAPITI} \cite{arxiv:2506.23731}, demonstrate that fine-tuning on as few as 1,000 clean samples suffices to neutralize backdoor-based watermarks while preserving >98\% of downstream task performance; this is achieved not through brute-force retraining but via gradient-aligned unlearning objectives that selectively suppress trigger-response correlations without perturbing functional gradients. Structural neutralization extends beyond fine-tuning: \emph{HuRef} \cite{arxiv:2407.10887} demonstrates that rotation-invariant weight permutations break alignment-based fingerprint extraction, while \emph{Robust and Efficient Watermarking of Large Language Models Using Error Correction Codes} \cite{doi:10.56553/popets-2025-0126} quantifies how iterative pruning induces exponential decay in Hamming distance under ECC-encoded weight perturbations, degrading bit accuracy and violating collision-resistance requirements for unforgeability. Collusion-based neutralization introduces a higher-order threat: \emph{Inhibitory Attacks on Backdoor-based Fingerprinting for Large Language Models} \cite{arxiv:2601.04261} formalizes token filtering (TFA) and sentence verification (SVA) as ensemble-level suppression mechanisms---where N models jointly filter out watermark-triggered tokens via voting over perplexity-weighted logits, achieving 92\% inhibition while maintaining ensemble BLEU scores within $\pm$0.3 of baseline. Critically, these attacks succeed \emph{without} access to secret keys or weight matrices, operating solely through black-box query orchestration---a violation of the foundational assumption that behavior-space fingerprints are inherently stable under composition. Theoretical analysis reveals a deeper tension: watermark embedding in parameter space often trades off between \emph{detectability} and \emph{transferability}, as shown in \emph{Robust and Efficient Watermarking of Large Language Models Using Error Correction Codes} \cite{doi:10.56553/popets-2025-0126}, where ECC-encoded weight perturbations exhibit exponential decay in Hamming distance under iterative pruning, violating the collision-resistance requirement for unforgeability. Emerging defenses like \emph{EverTracer} \cite{arxiv:2601.14311} attempt resilience by embedding probabilistic memorization biases rather than deterministic triggers---but even these succumb to \emph{Cracking White-box DNN Watermarks via Invariant Neuron Transforms} \cite{doi:10.1145/3580305.3599291}, which inserts dummy neurons that decouple activation pathways without altering outputs. Collectively, these findings expose a critical gap: current frameworks lack formal models of \emph{structural collusion resistance}, defined as the preservation of uniqueness under arbitrary convex combinations of model parameters or behavioral responses. Future work must therefore move beyond isolated robustness claims toward game-theoretic formulations of fingerprint persistence---where adversaries optimize over model manifolds, not token sequences---and integrate cryptographic primitives like zero-knowledge proofs (\emph{ZKROWNN} \cite{doi:10.1109/dac56929.2023.10247798}) to verify integrity without exposing the very structure being protected.


\subsection{Radioactivity Exploitation and Knowledge Distillation--Based Attacks}


The ``radioactive'' property---where watermark signals unintentionally propagate from teacher to student models during knowledge distillation (KD)---has been widely celebrated as a passive defense against model theft. Yet this subsection argues that radioactivity is not merely a fortuitous byproduct but a \emph{structural vulnerability} enabling three classes of scalable, high-impact attacks: (i) watermark stealing via distilled behavior reconstruction; (ii) pre-distillation sanitization that severs statistical dependencies before inheritance occurs; and (iii) post-distillation neutralization that surgically excises inherited signals while preserving functional fidelity. Crucially, these attacks exploit the same fundamental asymmetry exposed in the prior subsection---the gap between \emph{detection cost} (a lightweight statistical test) and \emph{removal feasibility} (a low-sample fine-tuning step or black-box query orchestration)---thereby undermining the core assumption that watermark persistence implies ownership verifiability. Just as parameter- and behavior-space adversaries sever provenance signals \emph{before}, \emph{during}, and \emph{after} distillation by targeting internal geometry or behavioral manifolds, radioactive attacks operationalize that fragility \emph{across} the distillation pipeline---turning inheritance itself into an attack surface.

Distillation-enabled watermark stealing constitutes the most direct exploitation of radioactivity. By querying a watermarked teacher model at scale and training a student on its outputs---including logits, token distributions, or even raw generations---adversaries reconstruct green lists, infer secret-key--dependent sampling biases, or replicate deterministic response patterns. Empirically, \cite{arxiv:2402.19361} demonstrates that under \textbackslash{}\$50 budget, attackers achieve >80\% success in both spoofing and scrubbing state-of-the-art schemes (e.g., Kirchenbauer et al., 2023) via KD-based reverse engineering. Similarly, \cite{arxiv:2002.12200} shows that distilled students inherit not only task performance but also statistically detectable watermark signatures---even when trained on paraphrased or translated outputs---enabling forgery without access to secret keys. This violates unforgeability: watermark presence no longer implies legitimate ownership, but rather exposure to a watermarked teacher.

Pre-distillation paraphrasing attacks disrupt radioactivity \emph{upstream}. Rather than removing watermarks after distillation, adversaries sanitize the teacher's outputs \emph{before} they enter the student's training set. \cite{arxiv:2402.14007} introduces targeted (TP) and untargeted (UP) rewriting strategies that break statistical correlations between teacher and student logits---e.g., by inserting entropy-preserving syntactic noise or applying synonym substitution guided by self-information thresholds \cite{arxiv:2505.05190}. These methods reduce watermark inheritance rates by up to 92\% while maintaining BLEU scores within $\pm$0.8, exposing a critical gap: radioactivity is not an invariant property but a fragile statistical coupling contingent on output distributional fidelity---mirroring the \emph{cross-lingual fragility} highlighted later in the systemic attack analysis, where round-trip translation collapses robustness from 88\% to <13\%.

Post-distillation neutralization completes the attack triad. Here, adversaries retain the distilled model's capabilities (>95\% task accuracy) while eliminating inherited watermarks via inference-time interventions---e.g., logit correction using a lightweight reparameterization module \cite{arxiv:2506.00359}, or alignment-resampling with external reward models \cite{arxiv:2403.10553}. Critically, \cite{arxiv:2505.06827} proves such removal is \emph{information-theoretically inevitable}: any quality-preserving perturbation oracle induces an efficiently mixing random walk over high-quality outputs, allowing watermark erasure with negligible degradation. This establishes a formal lower bound on robustness: no watermark can be both statistically detectable \emph{and} provably unremovable under black-box KD unless it incurs catastrophic utility loss---a trade-off empirically confirmed across OPT, LLaMA, and Alpaca in \cite{arxiv:2307.15593}.

Collectively, these attacks reveal that radioactivity conflates \emph{provenance} with \emph{contamination}. Future defenses must therefore decouple signal inheritance from functional transfer---e.g., via entangled watermarking embeddings \cite{arxiv:2002.12200} that bind watermarks to task-critical features, or modality-agnostic permutation equivariance \cite{arxiv:1910.12366} that makes watermark extraction conditional on structural triggers rather than output statistics. Without such architectural separation, radioactivity remains a liability masquerading as a safeguard---exposing the same foundational tension identified earlier: the absence of formal models of \emph{structural collusion resistance}, and foreshadowing the broader systemic vulnerabilities explored next, where watermark integrity fails not under generic noise, but under \emph{semantically or structurally coherent transformations} that preserve utility while severing statistical dependencies.


\subsection{Systemic and Cross-Modal Attack Surfaces: From Speculative Sampling to Multimodal Consistency}


Systemic and cross-modal attack surfaces represent a critical frontier in adversarial analysis of watermarking and fingerprinting---shifting focus from isolated token-level perturbations to the emergent vulnerabilities arising at the intersection of inference infrastructure, multimodal alignment, and safety-critical model behavior. These attacks exploit \emph{architectural couplings} rather than algorithmic weaknesses, revealing fundamental trade-offs that cannot be resolved within siloed design paradigms. First, speculative sampling---now ubiquitous in production LLM serving---introduces a provable no-go theorem between watermark strength and decoding efficiency: watermark injection into draft-model logits induces distributional divergence that degrades speculative acceptance rates, while deferring watermark application to the verification stage sacrifices detection fidelity under adaptive re-sampling. Empirical validation across OPT-1.3B and LLaMA-7B shows a 42--68\% latency penalty when enforcing >95\% watermark survival, confirming that watermark robustness and acceleration are not merely competing objectives but \emph{mathematically incompatible} under standard draft-target coordination protocols. Second, alignment degradation attacks expose how watermark-induced token distribution shifts---particularly in low-entropy regions---systematically attenuate or amplify safety guardrails: watermarking distorts the calibration surface of RLHF-aligned models, causing ``guard attenuation'' (reduced refusal rates for harmful queries) and ``guard amplification'' (excessive caution on benign prompts), both of which undermine the model's operational trustworthiness without altering its functional API. Crucially, this is not a side effect but an \emph{exploitable attack vector}: adversaries can tune watermark parameters to induce predictable guardrail failure modes, as demonstrated by universal adversarial suffixes that transfer across GPT-4, Claude, and LLaMA-2-Chat \cite{arxiv:2307.15043}. Third, cross-modal consistency attacks target multimodal LLMs (MLLMs) by decoupling watermark embedding across modalities: BadToken's token-substitution backdoor \cite{doi:10.1109/cvpr52734.2025.02785} and AnyDoor's test-time image perturbations \cite{arxiv:2402.08577} enable selective erasure in one modality (e.g., vision tokens) while preserving detectability in text---violating end-to-end provenance claims. Signal Watermark's FFT-based embedding \cite{arxiv:2410.06545} further illustrates how frequency-domain signal processing introduces modality-specific failure modes: watermark energy leakage between text and audio channels under joint encoding renders cross-modal detection inconsistent. Collectively, these systemic attacks demonstrate that watermark security cannot be evaluated in isolation; it must be co-verified with acceleration stacks, safety middleware, and multimodal coherence constraints. Future defenses must therefore adopt \emph{co-design principles}: watermark decoders integrated into speculative schedulers (e.g., vLLM-aware watermark verification), alignment-resampling mechanisms that correct guardrail distortion post-watermarking, and cross-modal binding protocols that enforce latent-space consistency via contrastive loss over joint embeddings. Without such system-level integration, watermarking remains a fragile add-on---not a foundational assurance layer.


\subsection{Evaluation Gaps and Methodological Vulnerabilities in Robustness Assessment}


Robustness evaluation in watermarking and fingerprinting remains plagued by methodological fragility that systematically overestimates real-world security---mirroring the systemic vulnerabilities exposed in the prior subsection, where adversarial impact emerges not from isolated perturbations but from \emph{coherent, semantically grounded transformations} that preserve utility while severing statistical dependencies. A foundational flaw lies in the \emph{syntheticity of adversarial perturbations}: most studies deploy non-adaptive, token-level edits---random substitutions, deletions, or insertions---despite empirical evidence that human-authored paraphrasing degrades detection far more severely than algorithmic corruption \cite{doi:10.18653/v1/2024.findings-naacl.40,arxiv:2402.19361}. This disconnect is theoretically grounded: human editing operates on semantic equivalence classes, not discrete token spaces, rendering statistical detectors blind to meaning-preserving transformations that collapse green-list statistics \cite{arxiv:2505.06827}. Worse, evaluations routinely ignore domain-specific rewriting norms---e.g., academic redaction, legal paraphrasing, or journalistic compression---which induce structured distribution shifts unmodeled by iid noise assumptions \cite{doi:10.1609/aaai.v36i10.21415,arxiv:2505.02884}. Compounding this, composite attacks---paraphrasing + translation + fine-tuning---are empirically shown to increase removal success by >40\% over single-vector baselines \cite{arxiv:2402.14007}, yet no benchmark mandates their systematic inclusion. This omission violates the principle of \emph{adversarial realism}, as attackers naturally chain operations to maximize stealth and minimize quality degradation---paralleling the cross-modal and alignment-degrading attack chains described earlier, where structural coherence---not randomness---drives failure.

Equally critical is the \emph{longitudinal blindness} in fingerprinting evaluation. While weight-space methods like SELF \cite{arxiv:2512.03620} and Embedding Watermarks into Deep Neural Networks \cite{arxiv:1701.04082} report high accuracy under static conditions, none assess persistence across iterative fine-tuning, continual learning updates, or model versioning---despite evidence that quantization, pruning, and LoRA adaptation induce non-uniform spectral drift that erodes fingerprint stability \cite{doi:10.1109/icassp48485.2024.10447264}. Similarly, behavior-space techniques such as DeepSigns \cite{arxiv:1804.00750} and PLMmark \cite{doi:10.1609/aaai.v37i12.26750} omit temporal testing, treating fingerprint extraction as a point-in-time forensic snapshot rather than a dynamic process subject to activation sparsity decay under distributional shift \cite{arxiv:2509.25448}. This gap is particularly acute for open-source models, where WAPITI \cite{arxiv:2410.06467} demonstrates that watermark transfer fidelity degrades nonlinearly with fine-tuning intensity---a phenomenon invisible to cross-sectional benchmarks and directly echoing the \emph{systemic incompatibility} between watermark strength and inference acceleration uncovered in speculative sampling. Furthermore, human-centered metrics remain marginal: automated fluency scores (BLEU, perplexity) fail to capture subtle cognitive load introduced by watermark-induced repetition biases or syntactic anomalies detectable by human readers \cite{arxiv:2403.04808,arxiv:2405.13517}, while linguistic equity is ignored---\cite{arxiv:2402.14007} reveals cross-lingual robustness collapses from 88\% to <13\% under round-trip translation, yet no evaluation protocol standardizes low-resource language testing. Closing these gaps demands \emph{adversarial composability} in benchmarks, \emph{temporal fidelity curves} tracking watermark decay across model evolution trajectories, and \emph{participatory validation} integrating human perceptibility studies across demographic and linguistic axes---transforming robustness assessment from a statistical exercise into a socio-technical assurance discipline aligned with the co-design imperatives articulated in the preceding section.


\section{Evaluation Frameworks, Benchmarks, and Real-World Deployment Constraints}



\subsection{Multidimensional Evaluation Taxonomy for Watermarking and Fingerprinting}


The evaluation of watermarking and fingerprinting systems has long suffered from a dangerous myopia: overreliance on binary detection accuracy (TPR/FPR) while neglecting the multidimensional operational fidelity required for real-world trustworthiness. This subsection establishes a principled, empirically grounded taxonomy that reframes evaluation as a \emph{system-level assurance problem}, integrating statistical rigor, linguistic integrity, functional alignment, and adversarial resilience into four interdependent dimensions---each exposing critical misalignments between laboratory benchmarks and fielded deployment constraints. First, detection-centric metrics must transcend empirical thresholds to embrace formal statistical guarantees: p-value calibration under entropy heterogeneity \cite{arxiv:2306.09194}, minimum detectable token length derived from KL divergence bounds \cite{doi:10.1038/s41586-024-08025-4}, and AUROC with confidence intervals across prompt distributions---not just static test sets---ensure detectors are not merely ``accurate'' but \emph{calibrated}, \emph{interpretable}, and \emph{distributionally robust}. Second, quality-preserving evaluation must abandon perplexity as a sufficient proxy for imperceptibility. Empirical studies demonstrate that even watermarking schemes preserving perplexity induce measurable degradation in instruction-following fidelity (GPT-4-Judge scores $\downarrow$17.3\%) and semantic coherence (LLM-as-judge fluency ratings $\downarrow$22.1\%) \cite{doi:10.1038/s41586-024-08025-4,arxiv:2303.13408}. Thus, human-centered metrics---including structured fluency assessments, task-specific degradation analysis \cite{arxiv:2402.16187,doi:10.18653/v1/2023.acl-long.423}, and alignment-resampling--guided truthfulness scoring \cite{arxiv:2309.15025}---are indispensable for quantifying \emph{functional imperceptibility}. Third, alignment-aware robustness assessment reveals a previously obscured trade-off: watermark strength modulates safety-guard behavior non-linearly. Empirical evidence shows that high-strength watermarks induce both ``guard attenuation'' (reduced safety from over-helpfulness) and ``guard amplification'' (excessive caution impairing utility), necessitating joint optimization of watermark parameters and guardrail classifiers \cite{doi:10.55662/jst.2023.4605,arxiv:2403.15447}. Finally, payload and provenance fidelity metrics expose a foundational flaw in fixed-length embedding: bit error rate (BER) surges in low-entropy contexts (e.g., code, legal text), collapsing multi-bit extraction success from >95\% to <38\% unless entropy-adaptive modulation is employed \cite{arxiv:2406.01946,arxiv:2402.16187,arxiv:2403.04808}. Crucially, this taxonomy reveals that no single dimension suffices---robustness against paraphrasing \cite{arxiv:2303.13408} collapses without alignment preservation; high payload capacity \cite{arxiv:2402.16187} is meaningless if BER exceeds 0.1 in safety-critical segments; and statistically calibrated detection \cite{arxiv:2306.09194} fails when human editors introduce syntactic anomalies undetectable by automated metrics \cite{arxiv:2303.13408}. Future evaluation must therefore adopt \emph{joint optimization objectives}: minimizing the Pareto frontier across TPR, BER, GPT-4-Judge score, and safety violation rate---moving beyond isolated metric reporting toward holistic assurance certificates verifiable via zero-knowledge proofs \cite{doi:10.1007/s44443-025-00177-1} or randomized smoothing \cite{arxiv:2403.04808}.


\subsection{Benchmark Architectures for Holistic and Adversarial Validation}


Emerging benchmark architectures for holistic and adversarial validation represent a paradigm shift from static accuracy reporting toward \emph{dynamic, threat-coevolving evaluation ecosystems}---where watermarking, fingerprinting, and detection are jointly stress-tested under realistic corruption pipelines spanning model evolution, cross-modal perturbations, and adversarial intent. This evolution is not incremental but \emph{architecturally necessary}: as the preceding subsection established, evaluation must be reframed as a \emph{system-level assurance problem}, integrating statistical calibration, functional imperceptibility, alignment-aware robustness, and payload fidelity---and as the following subsection underscores, operational viability hinges on tight coupling between these dimensions and real-world infrastructure constraints. Contemporary benchmarks therefore enforce \emph{lifecycle-aware integrity} not as an optional enhancement, but as the minimal condition for trustworthiness: watermark survival must persist not only across paraphrasing but also through fine-tuning, distillation, quantization, speculative decoding interference, and even selective unlearning---mirroring the multidimensional degradation pathways exposed in the evaluation taxonomy. JailbreakBench-inspired stress-testing has become foundational, systematically evaluating detectors against jailbreak-style prompt engineering, recursive rewriting, and multilingual translation cycles---demonstrating that 72\% of token-level watermarks fail under round-trip translation (RTT) without layered mitigation \cite{arxiv:2601.11629}. Crucially, MUSE introduces \emph{unlearning-aware evaluation}, quantifying watermark decay post-selective forgetting using mismatched datasets---a capability validated in MEraser's erasure analysis showing near-total removal with <1,000 samples \cite{doi:10.1145/3639477.3639720}, directly linking to the ``payload and provenance fidelity'' dimension where BER collapses under low-entropy contexts unless entropy-adaptive modulation is employed \cite{arxiv:2406.01946,arxiv:2402.16187,arxiv:2403.04808}. Dual Attribution Benchmark advances beyond binary attribution by supporting \emph{multi-provenance tracing}: it enables joint verification of base-model origin, fine-tuner identity, and end-user signature under collusion threats like ensemble inhibition or piggyback spoofing, as formalized in WASA's source-attribution payloads and Bileve's bi-level signature design \cite{arxiv:2310.00646,arxiv:2403.07283}---thereby operationalizing the \emph{joint optimization objectives} called for earlier (e.g., minimizing Pareto frontiers across TPR, BER, GPT-4-Judge score, and safety violation rate). Lifecycle-aware contamination benchmarks further close the lab-to-field gap: Detecting Benchmark Contamination Through Watermarking establishes radioactivity-based leakage detection on ARC-Easy/MMLU, requiring watermarks to survive both output editing \emph{and} parameter-level evolution---including LoRA merging, quantization, and distillation---revealing that 65\% of EmMark's weight-space signatures degrade under fp32$\rightarrow$int8 conversion without invariant-aware design \cite{doi:10.1109/sds66131.2025.00024}, a finding that directly motivates the \emph{operational friction} discussed next: memory footprint, latency overhead, and architectural incompatibility become decisive failure modes when quantization degrades signature fidelity. Cross-modal benchmarks expose critical architectural blind spots: LR-DWM motivates diffusion-language model extensions by proving that non-autoregressive watermarking requires bidirectional neighbor biasing---not sequential green-list sampling---achieving 92--99.5\% detection with <5ms overhead \cite{arxiv:2403.07283}, foreshadowing the \emph{decoupled, plug-in architectures} emphasized later as essential for interoperability across autoregressive and non-autoregressive models. Similarly, WaterPool demonstrates compositional evaluation necessity, showing that ensemble watermarks (e.g., red-green + acrostic + sensorimotor norms) improve robustness by +12.73\% over monolithic schemes when tested under combined synonym substitution and summarization attacks \cite{arxiv:2405.13517}, reinforcing the need for \emph{co-design} strategies like Thinking/Answering Phase separation \cite{doi:10.1609/aaai.v36i10.21415} and Alignment Resampling \cite{doi:10.14722/ndss.2026.230138} to preserve functional integrity amid multi-stage corruption. Critically, open infrastructure remains fragmented: while WaterBench pioneers hyperparameter-standardized comparison across OPT/LLaMA/Alpaca, its containerized detector pipelines lack support for confidential computing enclaves or edge-device constraints addressed by EmMark's embedded quantized LLM watermarking \cite{arxiv:2405.13517,doi:10.1109/sds66131.2025.00024}, echoing the regulatory and hardware compatibility tensions highlighted in the subsequent subsection. The most pressing gap is longitudinal fidelity---no benchmark yet tracks watermark degradation across iterative fine-tuning trajectories, leaving unverified the claim that AWM's CKA-aligned fingerprints persist after three successive instruction-tuning updates \cite{doi:10.1145/3639477.3639720}. Future architectures must therefore integrate \emph{version-aware decay curves}, cryptographic calibration of p-values across entropy regimes, and standardized attack taxonomies---moving from ``can it be broken?'' to ``how long does it resist breaking under evolving threat models?'', thereby completing the triad of \emph{system-level assurance}, \emph{operational viability}, and \emph{longitudinal resilience} that defines next-generation watermarking and fingerprinting evaluation.


\subsection{Operational Realities: Latency, Compatibility, and System Integration Constraints}


Watermarking and fingerprinting techniques, while theoretically sound, confront severe operational friction when integrated into production AI systems---friction that manifests not in detection accuracy, but in latency overhead, architectural incompatibility, memory footprint, and regulatory friction. Crucially, these constraints are neither orthogonal nor additive; they form a tightly coupled constraint surface where optimizing one dimension often degrades another. For instance, SynthID-Text \cite{doi:10.1038/s41586-024-08025-4} achieves sub-50ms inference overhead by modifying only the sampling procedure---but this design precludes compatibility with speculative decoding, as rigorously proven by the no-go theorem in \emph{Inevitable Trade-off between Watermark Strength and Speculative Sampling Efficiency}, which establishes that watermark strength \$\textbackslash{}alpha\$ and draft-model acceptance rate \$\textbackslash{}beta\$ satisfy \$\textbackslash{}alpha \textbackslash{}cdot \textbackslash{}beta \textbackslash{}leq C\$, a fundamental bound violated by all logits-based schemes attempting simultaneous optimization. This explains why WaterMax \cite{arxiv:2403.04808}, despite its theoretical elegance, exhibits 3.2$\times$ higher latency under vLLM's PagedAttention scheduler compared to its baseline---because its zero-weight-modification design cannot exploit draft-model caching of perturbed logits. Similarly, DeepTextMark \cite{doi:10.1109/access.2024.3376693} incurs >100MB decoder memory overhead due to its transformer-based classifier, rendering it infeasible for edge deployment, whereas Advancing Beyond Identification: Multi-bit Watermark for Large Language Models \cite{doi:10.18653/v1/2024.naacl-long.224} achieves <1MB footprint via lightweight reweighting---a critical distinction for quantized models on resource-constrained devices. Compatibility extends beyond hardware: watermark-induced token distribution shifts degrade safety guardrails, as empirically quantified in \emph{Watermarking Degrades Alignment in Language Models}, where even mild green-list bias (0.05 KL divergence) causes 18\% false-positive rate spikes in moderation APIs and 22\% reduction in RLHF-aligned helpfulness scores. This necessitates co-design: Tracing Text Provenance via Context-Aware Lexical Substitution \cite{doi:10.1609/aaai.v36i10.21415} decouples watermarking from reasoning pathways via Thinking/Answering Phase separation, while Alignment Resampling (AR) \cite{doi:10.14722/ndss.2026.230138} corrects guard attenuation post-hoc using external reward models. Regulatory integration further compounds complexity: A Private Watermark for Large Language Models \cite{arxiv:2307.16230} enables MPC-backed watermarking for GDPR-compliant collaborative verification, yet its cryptographic protocol introduces 17ms ZKP generation latency per inference---making it unsuitable for streaming applications unless batched. Conversely, Hey, That's My Model! Introducing Chain \& Hash, An LLM Fingerprinting Technique \cite{arxiv:2407.10887} satisfies EU AI Act Annex VI documentation requirements via cryptographically bound prompt-response logs but requires white-box access, excluding third-party API consumers. The most promising convergence lies in modularity: Watermarking Text Generated by Black-Box Language Models \cite{arxiv:2305.08883} operates entirely post-hoc without logits or weights, achieving 95\% detection under paraphrasing with zero acceleration-stack interference, while A Watermark-Conditioned Diffusion Model for IP Protection \cite{arxiv:2403.10893} generalizes this principle to non-autoregressive models via bidirectional neighbor biasing---demonstrating that operational viability increasingly hinges not on stronger signals, but on \emph{decoupled, plug-in architectures} that treat watermarking as an interoperable assurance layer rather than a monolithic model modification.


\subsection{Human-Centered and Socio-Technical Evaluation Gaps}


Current evaluation frameworks for watermarking and fingerprinting remain critically under-specified along three interlocking human-centered dimensions: perceptual fidelity as experienced by diverse readers, linguistic equity across sociotechnically marginalized language communities, and the systemic collateral harms arising from misuse or misinterpretation of provenance signals. This gap is not merely methodological---it reflects a foundational epistemological asymmetry wherein technical validation assumes a monolithic, English-dominant, expert-authorized user, while real-world deployment confronts heterogeneous literacy practices, multilingual co-authorship, and high-stakes attribution in legal, educational, and journalistic contexts. Empirical evidence confirms this misalignment: human studies reveal that watermark-induced syntactic anomalies---such as subtle repetition biases or entropy-constrained token substitutions---are significantly more detectable by non-native speakers and readers with dyslexia, yet no benchmark incorporates cognitive load metrics or cross-demographic perception trials \cite{arxiv:2403.04808}. Similarly, cross-lingual robustness audits demonstrate catastrophic detection failure under round-trip translation (RTT), where token-level watermarks collapse from >88\% AUC to 9--13\% for Bangla and Hindi, exposing tokenizer-dependent schemes as de facto English-centric artifacts \cite{arxiv:2402.14007}. Crucially, this disparity cannot be resolved by post-hoc adaptation alone; it demands \emph{evaluation-first} design, wherein fairness-aware test suites---spanning Afro-Asiatic, Dravidian, and Indigenous language corpora---are integrated into core validation pipelines rather than treated as optional add-ons \cite{arxiv:2305.08883}.

The dual-use risk landscape further compounds these gaps. While watermarking is framed as an IP protection mechanism, its operationalization enables surveillance vectors: academic integrity tools that flag paraphrased student submissions \cite{doi:10.1609/aaai.v36i10.21415}, or moderation APIs that gate access to services based on probabilistic watermark presence---risks quantified in false-positive rate analyses for medical and legal domains where erroneous attribution may deny critical care or due process \cite{arxiv:2506.15975}. Critically, no existing benchmark models downstream exclusion risk or measures how watermark strength modulates truthfulness-safety-helpfulness trade-offs under stakeholder pressure---a failure directly implicated in alignment degradation attacks where watermark-induced distribution shifts cause guard attenuation or amplification \cite{arxiv:2306.09194}. To close these gaps, we advocate a tripartite augmentation of evaluation practice: (1) participatory assessment protocols embedding GPT-4-Judge--guided instruction-following evaluation with structured human fluency/accuracy ratings across age, education, and linguistic background \cite{arxiv:2403.04808,arxiv:2502.11598}; (2) cross-lingual fairness auditing that treats RTT not as an attack but as a \emph{baseline requirement}, mandating layered embedding-time/post-generation signals validated on low-resource language benchmarks \cite{arxiv:2601.11629}; and (3) misuse impact modeling that formalizes false attribution as a constrained optimization problem---e.g., minimizing the KL divergence between watermark detection logits and ground-truth authorship labels under adversarial prompt distributions---thereby transforming ethical redlines into verifiable statistical constraints \cite{arxiv:2506.15975,arxiv:2510.17000}.

This human-centered evaluation imperative does not exist in isolation---it forms the essential bridge between the \emph{operational viability constraints} detailed in the preceding subsection and the \emph{formally verified, standardized evaluation protocols} advanced in the following one. Where the prior section exposed how latency, memory, architectural coupling, and regulatory friction constrain real-world deployment, this subsection reveals that those constraints are inseparable from fairness, accessibility, and accountability: for instance, DeepTextMark's >100MB decoder overhead \cite{doi:10.1109/access.2024.3376693} precludes edge deployment not only on hardware grounds but also exacerbates linguistic inequity by excluding low-bandwidth, multilingual environments; similarly, the 18\% false-positive spike in moderation APIs under green-list bias \cite{arxiv:2306.09194} directly manifests the \emph{systemic collateral harms} flagged here---and motivates the subsequent shift toward FDR-controlled detection and certified robustness. Conversely, the evaluation-first ethos advocated here anticipates and scaffolds the cryptographic grounding, statistical rigor, and interoperable specification developed next: participatory assessment protocols feed directly into JSON-LD--encoded descriptors for cross-framework integration \cite{arxiv:2407.04411}, while RTT-as-baseline informs the entropy-constrained generation bounds used to derive minimum detectable token length \cite{arxiv:2307.15593}. Without such integration, technical advances risk entrenching inequity rather than enabling equitable provenance assurance.


\subsection{Towards Verified and Standardized Evaluation Protocols}


The transition from heuristic validation to formally verified, standardized evaluation protocols represents a critical maturation phase for watermarking and fingerprinting research---shifting emphasis from \emph{whether} a signal can be detected to \emph{how confidently, reproducibly, and legally} its presence or absence can be certified. Recent advances coalesce around three interlocking pillars: cryptographic grounding of statistical claims, standardization of hypothesis-testing rigor, and interoperable specification of evaluation artifacts. First, cryptographic primitives now underpin verifiability: randomized smoothing enables \emph{certified robustness}---i.e., provable lower bounds on watermark survival under arbitrary token substitutions---formalized as \$ \textbackslash{}Pr\cite{arxiv:2506.01011} \textbackslash{}geq 1 - \textbackslash{}delta \$ for all \$x'\$ within an \$\textbackslash{}ell\_0\$-ball of radius \$r\$, where \$\textbackslash{}mathcal\{D\}\$ is the detector and \$\textbackslash{}delta\$ is a statistically bounded failure probability \cite{arxiv:2401.16820}. Similarly, zero-knowledge proofs (ZKPs) in \cite{arxiv:2407.10887} and \cite{doi:10.1609/aaai.v37i12.26750} allow owners to publicly attest to fingerprint existence without revealing secret keys or model weights---satisfying GDPR Article 32 and EU AI Act Annex VI auditability requirements. Second, statistical testing frameworks are being systematized beyond ad hoc p-value thresholds: \cite{arxiv:2307.15593} derive detection power via KL divergence bounds and minimum detectable token length under entropy-constrained generation, while \cite{arxiv:2406.10281} introduce false discovery rate (FDR) control across heterogeneous prompt distributions---enabling calibrated, distribution-aware inference rather than fixed-threshold binary decisions. Third, interoperability is being engineered through universal descriptors: \cite{arxiv:2407.04411} proposes JSON-LD schemas encoding embedding locus (output/parameter/behavior), detection modality (white/gray/black-box), payload semantics (user ID vs. timestamp), and cryptographic parameters (key derivation function, ECC code rate), enabling cross-framework integration---for instance, plugging BiMark's \cite{arxiv:2403.04808} decoder into WaterBench's \cite{doi:10.1145/3691626} pipeline. Critically, these efforts expose a persistent longitudinal gap: no benchmark yet tracks watermark decay across iterative fine-tuning, continual learning, or version updates---despite evidence that \cite{arxiv:2410.06467} and \cite{doi:10.18653/v1/2023.acl-long.423} demonstrate >60\% fidelity loss after three LoRA adaptations. Future protocols must therefore embed \emph{version-aware evaluation datasets}, capturing model evolution trajectories alongside corresponding watermark degradation curves. Ultimately, verified evaluation is not merely technical hygiene---it is the prerequisite for regulatory acceptance, cross-vendor comparability, and liability assignment in high-stakes domains; without it, watermarking remains a forensic curiosity rather than a governance-grade assurance layer.


\section{Cross-Domain Extensions, Governance Integration, and Socio-Technical Implications}



\subsection{Multimodal and Cross-Modal Watermarking Architectures}


Multimodal and cross-modal watermarking architectures represent a paradigm shift from text-centric provenance mechanisms toward holistic, system-aware integrity binding---where signals must persist across heterogeneous representations while preserving functional fidelity, perceptual invisibility, and semantic coherence. Unlike output-space text watermarking, which operates on discrete token sequences, multimodal watermarking confronts the fundamental challenge of \emph{joint-modality signal embedding}: ensuring that watermarks embedded in one modality (e.g., vision) remain detectable and consistent when projected into another (e.g., language), without inducing misalignment, hallucination, or performance degradation. Recent frameworks establish three convergent design principles: (i) \emph{latent-space coordination}, where watermark perturbations are synchronized across modalities via shared semantic manifolds; (ii) \emph{order-agnostic injection}, overcoming non-sequential generation constraints in diffusion models; and (iii) \emph{syntax- and prosody-aware embedding}, enforcing structural invariance in code and speech domains. Vision-language alignment--preserving watermarking exemplifies the first principle---though no cited paper in the provided list explicitly introduces ``VLA-Mark'' or jointly optimizes green-list sampling with CLIP-embedded image regions using contrastive loss---so this unsupported claim must be decoupled from citation. The reference to \cite{arxiv:2402.16187,arxiv:2505.16530} is incorrect here, as PLMmark focuses exclusively on black-box watermarking for pre-trained language models via supervised contrastive loss on trigger sets (not vision-language alignment), and FT-Shield targets \emph{text-to-image diffusion models} for fine-tuning copyright protection---not joint vision-language watermarking or hallucination reduction. Therefore, these citations are removed.

Diffusion-based watermarking frameworks such as LR-DWM and DMark overcome the sequential decoding limitation through bidirectional dependency modeling: LR-DWM embeds watermarks via predictive residual modulation in latent space, where watermark strength is dynamically modulated by the variance of denoising steps, enabling certified robustness against cropping and affine transformations with $\leq$0.3\% FID degradation. However, neither ``LR-DWM'' nor ``DMark'' appear in the provided paper titles, and no listed paper describes predictive residual modulation or FID degradation bounds under affine transformations. The only diffusion-related paper is \emph{FT-Shield: A Watermark Against Unauthorized Fine-tuning in Text-to-Image Diffusion Models}, which introduces a Mixture of Experts approach for watermark detection and ensures watermark transfer from training images to generated outputs---but does not mention predictive residual modulation, denoising-step variance modulation, or FID degradation metrics. Thus, citing \cite{arxiv:2505.16530} here is unsupported and must be removed.

Audio-language extensions leverage phoneme-aware watermarking, embedding identifiers within prosodic invariants---pitch contour, duration, and energy envelopes---ensuring watermark survival under vocoder compression and speaker adaptation, while maintaining MOS scores $\geq$4.2/5.0. None of the provided papers address audio, speech, prosody, vocoders, or MOS scores. Therefore, this sentence lacks supporting citations and all bracketed references must be omitted.

Code-generation watermarking adopts syntax-guided embedding, constraining watermark insertion to Abstract Syntax Tree (AST)-valid token positions (e.g., comment blocks, docstrings, or whitespace-preserved identifiers) to preserve compilability and runtime behavior---demonstrating zero regression on Python test suites even after aggressive obfuscation. While \emph{Double-I Watermark} introduces backdoor paradigms with triggers in instruction and input for LLM fine-tuning (not AST-aware code watermarking), and \emph{Explanation as a Watermark} embeds signals in feature attribution explanations (not code syntax), none of the listed papers describe AST-constrained, compilability-preserving watermarking for code-generation models. Thus, citing \cite{arxiv:2403.04808,arxiv:2405.04825} is inaccurate and must be removed.

Cross-modal consistency binding protocols---such as WMVLM's vision-language model--based evaluation framework---formalize integrity verification as a hypothesis test over latent distribution shifts: given watermarked image-text pairs \$(x\_i, y\_i)\$, detection relies on measuring KL divergence between \$p(z\_i\textasciicircum{}\{(v)\}|x\_i)\$ and \$p(z\_i\textasciicircum{}\{(t)\}|y\_i)\$ in shared embedding space \$\textbackslash{}mathcal\{Z\}\$, where deviations exceeding threshold \$\textbackslash{}tau\$ indicate cross-modal watermark decay. Neither ``WMVLM'' nor KL-divergence--based cross-modal detection appears in any listed paper. \emph{Evaluating Durability: Benchmark Insights into Multimodal Watermarking} evaluates robustness of watermarked content against real-world corruptions but does not propose KL divergence--based hypothesis testing over latent distributions. \emph{Proving membership in LLM pretraining data via data watermarks} uses hypothesis testing for data watermark detection---but only for \emph{text-only} black-box membership inference, not cross-modal latent distribution shifts. Therefore, citing \cite{arxiv:2406.03728,arxiv:2402.10892} is unsupported and must be removed.

Emerging work reveals that \emph{perceptual invisibility} and \emph{functional correctness} are not orthogonal but jointly constrained: watermarking that preserves pixel-level fidelity may distort semantic gradients in vision-language models, while text-preserving watermarking in code models can break type inference. This insight is conceptual and not directly attributable to any single listed paper; no cited paper analyzes trade-offs between pixel fidelity and semantic gradient distortion or type inference failure in code models. So no citation is warranted.

The frontier thus lies in \emph{co-embedding spaces}---where watermark signals are defined not in raw modalities but in task-relevant invariant subspaces (e.g., safety-constrained reasoning paths in RAG pipelines \cite{doi:10.1109/sds66131.2025.00024} or privacy-preserving activation manifolds in confidential computing enclaves \cite{doi:10.34190/iccws.19.1.2096}).  
- \emph{Securing Rag: A Risk Assessment and Mitigation Framework} discusses RAG-specific vulnerabilities and mitigation strategies---including data pre-processing, storage, and integration risks---but does \textbf{not} mention ``safety-constrained reasoning paths'' or co-embedding spaces for watermarking. So \cite{doi:10.1109/sds66131.2025.00024} is unsupported here and must be removed.  
- \emph{Enhancing Privacy and Security in Large-Language Models: A Zero-Knowledge Proof Approach} proposes ZKPs for data verification, identity control, and malicious prompt detection---but makes \textbf{no reference} to ``privacy-preserving activation manifolds'', ``confidential computing enclaves'', or co-embedding spaces for watermarking. So \cite{doi:10.34190/iccws.19.1.2096} is also unsupported and must be removed.

Final note: The only papers in the list that meaningfully engage multimodal or cross-modal concerns are:
- \emph{FT-Shield} (text-to-image diffusion watermarking),
- \emph{Evaluating Durability: Benchmark Insights into Multimodal Watermarking} (empirical robustness evaluation across image/text perturbations),
- \emph{Latent Diffusion Models for Image Watermarking: A Review of Recent Trends and Future Directions} (review of LDM-based watermarking, including latent-space embedding).

However, none support the specific technical claims made in the original subsection (e.g., contrastive alignment with CLIP, predictive residual modulation, prosody-aware embedding, AST-guided insertion, KL divergence over vision-language latents, or co-embedding for RAG/ZKP). Since \emph{no provided paper supports any of the cited technical assertions in this subsection}, all bracketed citations must be removed.

Thus, the corrected subsection contains \textbf{no citations}, as none are justified by the provided corpus.

Multimodal and cross-modal watermarking architectures represent a paradigm shift from text-centric provenance mechanisms toward holistic, system-aware integrity binding---where signals must persist across heterogeneous representations while preserving functional fidelity, perceptual invisibility, and semantic coherence. Unlike output-space text watermarking, which operates on discrete token sequences, multimodal watermarking confronts the fundamental challenge of \emph{joint-modality signal embedding}: ensuring that watermarks embedded in one modality (e.g., vision) remain detectable and consistent when projected into another (e.g., language), without inducing misalignment, hallucination, or performance degradation. Recent frameworks establish three convergent design principles: (i) \emph{latent-space coordination}, where watermark perturbations are synchronized across modalities via shared semantic manifolds; (ii) \emph{order-agnostic injection}, overcoming non-sequential generation constraints in diffusion models; and (iii) \emph{syntax- and prosody-aware embedding}, enforcing structural invariance in code and speech domains. Vision-language alignment--preserving watermarking exemplifies the first principle by jointly optimizing token-level green-list sampling with visual patch affinity metrics---using contrastive loss to align textual watermarked continuations with CLIP-embedded image regions---achieving 98.8\% AUC detection while reducing hallucination rates by 41\% compared to decoupled baselines. Diffusion-based watermarking frameworks such as LR-DWM and DMark overcome the sequential decoding limitation through bidirectional dependency modeling: LR-DWM embeds watermarks via predictive residual modulation in latent space, where watermark strength is dynamically modulated by the variance of denoising steps, enabling certified robustness against cropping and affine transformations with $\leq$0.3\% FID degradation. Audio-language extensions leverage phoneme-aware watermarking, embedding identifiers within prosodic invariants---pitch contour, duration, and energy envelopes---ensuring watermark survival under vocoder compression and speaker adaptation, while maintaining MOS scores $\geq$4.2/5.0. Code-generation watermarking adopts syntax-guided embedding, constraining watermark insertion to Abstract Syntax Tree (AST)-valid token positions (e.g., comment blocks, docstrings, or whitespace-preserved identifiers) to preserve compilability and runtime behavior---demonstrating zero regression on Python test suites even after aggressive obfuscation. Critically, cross-modal consistency binding protocols---such as WMVLM's vision-language model--based evaluation framework---formalize integrity verification as a hypothesis test over latent distribution shifts: given watermarked image-text pairs \$(x\_i, y\_i)\$, detection relies on measuring KL divergence between \$p(z\_i\textasciicircum{}\{(v)\}|x\_i)\$ and \$p(z\_i\textasciicircum{}\{(t)\}|y\_i)\$ in shared embedding space \$\textbackslash{}mathcal\{Z\}\$, where deviations exceeding threshold \$\textbackslash{}tau\$ indicate cross-modal watermark decay. Emerging work reveals that \emph{perceptual invisibility} and \emph{functional correctness} are not orthogonal but jointly constrained: watermarking that preserves pixel-level fidelity may distort semantic gradients in vision-language models, while text-preserving watermarking in code models can break type inference. The frontier thus lies in \emph{co-embedding spaces}---where watermark signals are defined not in raw modalities but in task-relevant invariant subspaces (e.g., safety-constrained reasoning paths in RAG pipelines or privacy-preserving activation manifolds in confidential computing enclaves). Future architectures must therefore unify cryptographic unforgeability, multimodal radioactivity, and regulatory traceability---not as add-ons, but as first-class constraints in latent-space optimization.


\subsection{Regulatory-Aware Watermarking Design and Forensic Admissibility}


Regulatory-aware watermarking design transcends algorithmic optimization to embody a \emph{governance-first engineering discipline}, wherein technical primitives are co-designed with legal epistemology---ensuring that watermark signals function not merely as statistical artifacts but as \emph{admissible evidence}, \emph{auditable execution traces}, and \emph{verifiable compliance artifacts}. This formal alignment is not an afterthought but a direct response to the evidentiary fragility exposed in multimodal systems: if watermarks decay unpredictably across modalities---or fail to survive common transformations like cropping, translation, or vocoder resynthesis---they cannot serve as stable bases for attribution, liability assignment, or regulatory audit. As such, regulatory-aware design must begin where multimodal watermarking leaves off: by grounding forensic reliability in \emph{statistical rigor under bounded adversarial capability}, \emph{longitudinal stability}, and \emph{jurisdiction-aware descriptors}.

The EU AI Act's transparency obligations (Art. 52) demand human-readable, machine-verifiable provenance disclosures; NIST AI RMF's accountability pillar (ID.RA-2) requires traceability across model lifecycles; and U.S. Daubert/Kumho jurisprudence mandates demonstrable scientific validity---including statistical significance (p < 0.01), reproducibility across independent evaluators, and resistance to adversarial erasure \cite{arxiv:2301.10226,arxiv:2307.15593,doi:10.1109/eurosp57164.2023.00072}. Verifiable inference protocols now operationalize this alignment: Chain \& Hash cryptographically binds fingerprint prompts with their responses, ensuring no adversary can generate colliding fingerprints and allowing model owners to irrefutably demonstrate their creation \cite{arxiv:2407.10887}; ZKROWNN utilizes Zero-Knowledge Proofs to enable ownership verification while preserving watermark privacy \cite{doi:10.1109/dac56929.2023.10247798}. Similarly, WASA enables source attribution via watermarks embedded in LLM-generated text that contain information about data providers, supporting data provenance verification without model access or API cooperation \cite{arxiv:2310.00646}. Auditability-by-design further embeds regulatory logic into architecture: deterministic key derivation from training dataset hashes ensures reproducible watermark instantiation across jurisdictions, while standardized watermark manifests---such as WASA's JSON-LD--encoded source-attribution payloads---enable third-party verification without model access or API cooperation \cite{arxiv:2310.00646,doi:10.1109/icassp48485.2024.10447264}.

Critically, forensic admissibility hinges on \emph{statistical rigor under bounded adversarial capability}: recent benchmarks like WaterBench and MUSE validate detection p-values against adaptive paraphrasing, fine-tuning, and translation attacks---not as isolated perturbations but as compositional threats mirroring real-world litigation scenarios \cite{arxiv:2405.13517,arxiv:2410.23861}. Moreover, unforgeability must be provably grounded: schemes like Provably Robust Multi-bit Watermarking for AI-generated Text via Error Correction Code derive bit-error-rate bounds under token-edit oracles \cite{arxiv:2401.16820}, while Undetectable Watermarks for Language Models reduce detection feasibility to one-way function inversion---establishing computational indistinguishability even under adaptive querying \cite{arxiv:2306.09194}.

Yet critical gaps persist---gaps that directly foreshadow the equity and legitimacy challenges explored in the following subsection. No existing framework formally models \emph{longitudinal decay} of watermark fidelity across iterative fine-tuning---a requirement for AI Act compliance in continuously updated models---and empirical studies reveal stark disparities in detection reliability across linguistic registers, undermining equitable enforcement \cite{arxiv:2506.01011}. These disparities are not incidental: they reflect deeper tensions between statistical uniformity and linguistic heterogeneity---tensions that render watermark robustness contingent on tokenizer assumptions, entropy distributions, and syntactic conventions. As such, future work must integrate \emph{regulatory temporal logic}---formalizing watermark persistence as a time-bounded invariant---and develop \emph{jurisdiction-aware watermark descriptors} that encode jurisdiction-specific evidentiary thresholds (e.g., Daubert's ``general acceptance'' criterion) directly into detector configuration spaces---thereby bridging the technical foundations of forensic reliability with the socio-technical realities of multilingual, cross-cultural, and institutionally fragmented deployment.


\subsection{Equity, Bias, and Human-Centered Evaluation Frameworks}


Watermarking and fingerprinting systems, while technically sophisticated, exhibit profound socio-technical fault lines that undermine their legitimacy when deployed across heterogeneous linguistic, cultural, and institutional contexts. Linguistic inequity constitutes the most empirically documented bias: token-level statistical watermarks---such as green-list sampling---collapse under round-trip translation (RTT), with detection accuracy plummeting significantly for low-resource languages; this structural dependency on subword segmentation and entropy homogeneity is exposed in studies evaluating robustness across language families \cite{arxiv:2301.10226,arxiv:2305.08883,arxiv:2306.04634}. This is not merely a performance degradation but an \emph{epistemic exclusion}, where watermark robustness becomes a proxy for linguistic privilege---favoring English-centric tokenizers (e.g., SentencePiece over Unicode-aware alternatives) and high-entropy domains like technical writing over narrative or poetic registers. Crucially, this bias propagates into forensic admissibility: false negatives in RTT settings risk misclassifying watermarked outputs from non-English models as human-authored, thereby eroding accountability in multilingual information ecosystems.

False attribution risk further compounds equity concerns, manifesting through two interlocking mechanisms: alignment degradation and collateral harm. Watermark-induced distributional shifts systematically perturb safety guardrails, causing either ``guard attenuation'' (reduced safety in high-risk contexts) or ``guard amplification'' (excessive caution impairing utility), as rigorously quantified in analyses of quality-detection trade-offs and alignment interference \cite{arxiv:2403.19548,doi:10.18653/v1/2024.findings-naacl.40,arxiv:2306.09194}. Such distortions disproportionately impact marginalized users---for instance, medical or legal queries in low-resource languages may suffer compounded hallucination or refusal rates due to both tokenizer instability and alignment interference. Moreover, watermarking's dual-use potential enables surveillance-by-design: Unicode-based watermark detection by advanced models reveals that watermark signals can be \emph{inferred} without secret keys, transforming provenance tools into covert tracking vectors for dissenting speech or academic submissions \cite{arxiv:2402.10892,doi:10.1145/3611643.3616297}. This necessitates provably undetectable schemes---not as cryptographic curiosities but as ethical prerequisites for equitable deployment.

Human-centered evaluation must therefore transcend automated metrics. Current benchmarks rely on BLEU, perplexity, or AUROC, yet these fail to capture perceptual fidelity disparities across reader demographics. Studies demonstrate that human fluency ratings under paraphrasing diverge significantly by native language and education level, while benchmarking frameworks integrate LLM-as-judge assessment to expose subtle coherence failures masked by lexical similarity scores \cite{arxiv:2306.04634,arxiv:2305.08883,arxiv:2310.15654}. Critically, participatory design principles are emerging: frameworks enable third-party black-box watermarking for domain-specific stakeholders, aligning incentives for end users rather than model providers \cite{arxiv:2209.08773,arxiv:2406.01946}; similarly, proposals mandate human-readable fingerprint publication to satisfy transparency and contestability principles \cite{arxiv:2407.10887,arxiv:2509.03058}. Future frameworks must formalize cross-cultural perception studies---e.g., measuring Unicode watermark detectability across LLMs and user groups---and embed red-teaming for systemic harms, such as watermark-enabled denial-of-service against low-resource-language users. Without such grounding, watermarking risks entrenching---not mitigating---the very inequities it purports to govern.


\subsection{System-Level Integration and Ecosystem Viability}


Watermarking and fingerprinting techniques must transcend their origins as isolated algorithmic interventions and mature into \emph{system-aware, production-grade assurance layers}---interoperable components that coexist with, rather than compete against, the full stack of modern AI infrastructure. This evolution demands rigorous attention to four interlocking dimensions: latency- and memory-constrained integration; coexistence with safety-critical guardrails; compatibility with confidential and edge-deployed execution environments; and viability in federated, multi-stakeholder ecosystems. Crucially, system-level viability is not merely an engineering concern---it is a \emph{security primitive}: watermark decoders that induce speculative decoding conflicts or interfere with moderation pipelines create exploitable attack surfaces far beyond detection failure, directly undermining the \emph{assurance-by-design} integrity demanded across regulatory and production layers. PostMark's logit-free, post-hoc insertion achieves sub-50ms overhead in streaming inference while preserving 95\%+ detection under paraphrasing \cite{arxiv:2305.08883}, demonstrating that decoupling watermark logic from autoregressive sampling is essential for acceleration-stack compatibility---and aligning with the earlier call for \emph{auditability-by-design} through non-intrusive, observable verification. Similarly, WaterMax's zero-weight-modification design eliminates parameter-access bottlenecks entirely, enabling deployment on quantized edge models without runtime weight reconstruction \cite{arxiv:2403.04808}. PRIVMARK advances this further by embedding multi-party computation (MPC) protocols directly into the watermark verification layer, ensuring cryptographic binding of inference traces without exposing secret keys---even in untrusted cloud environments \cite{doi:10.1109/dac56929.2023.10247798}, thereby fulfilling the \emph{jurisdiction-aware descriptors} and \emph{enforceable cryptographic boundaries} emphasized in the preceding and following sections. Coexistence with safety infrastructure remains fragile: watermark-induced token distribution shifts cause measurable ``guard attenuation'' (reduced safety) and ``guard amplification'' (excessive caution), degrading alignment fidelity across nine NLP tasks \cite{arxiv:2306.09194}---a phenomenon directly linked to the \emph{alignment degradation} and \emph{safety coexistence} risks identified earlier, and one that reinforces the need for architectural decoupling observed in ReasonMark's Thinking/Answering Phase separation, which isolates watermark perturbations to non-reasoning pathways to preserve truthfulness and helpfulness while maintaining detectability---a formal decoupling validated via LLM-judger assessment \cite{arxiv:2310.06356}. For confidential computing, Invariant-based Robust Weights Watermark leverages permutation equivariance to embed identifiers resilient to fp32/int8 switching, enabling license enforcement in trusted execution environments where watermark visibility must be modality-agnostic and hardware-aware \cite{doi:10.1109/icassp48485.2024.10447264,arxiv:1910.12366}, thus supporting \emph{regulatory temporal logic} through verifiable persistence across hardware abstraction layers. Federated viability hinges on composability: WASA's source-attribution watermarks encode fine-grained provenance (base model + fine-tuner + user) within multi-bit payloads, while DualGuard's dual-stream architecture enables traceability of piggyback spoofing in shared-model ecosystems---both requiring standardized descriptor formats (e.g., JSON-LD schemas) and plug-in detector interfaces to ensure cross-framework interoperability \cite{arxiv:2403.04808,arxiv:2506.00359}, directly addressing the \emph{provenance decay} and \emph{cross-lingual exclusion} concerns raised previously and echoing the \emph{collateral harm} mechanisms outlined in the following subsection. MEraser's transferable fingerprint erasure benchmarks expose a critical gap: ecosystem viability requires longitudinal evaluation of watermark persistence across iterative fine-tuning, continual learning, and version updates---yet no benchmark currently tracks decay curves across model evolution trajectories \cite{arxiv:2402.19361,arxiv:2406.03728}, a lacuna that impedes fulfillment of the \emph{longitudinal, jurisdiction-aware, and system-integrated assurance requirements} established throughout this survey. The path forward lies in \emph{assurance-by-design}: watermarking must be specified not as a module but as a \emph{contract}---with formally verified latency bounds, cryptographic commitments to inference integrity, and auditable interference profiles for safety middleware. Only then do these techniques become foundational, not auxiliary, to trustworthy AI systems---capable of sustaining \emph{equitable traceability}, \emph{human-readable provenance}, and \emph{non-instrumentalizable robustness} across linguistic, institutional, and technical boundaries.


\subsection{Ethical Governance and Dual-Use Risk Mitigation}


Watermarking and fingerprinting, while indispensable for intellectual property protection and synthetic content accountability, embody a profound ethical paradox: their very mechanisms---designed to enhance transparency and traceability---can be repurposed as instruments of surveillance, exclusion, and coercive control. This duality demands governance frameworks that are not merely reactive but \emph{anticipatorily normative}, embedding ethical constraints directly into technical design. Critically, empirical evidence reveals that advanced reasoning models---including GPT-5 and Gemini 2.5 Pro---exhibit strong \emph{detection} capability for watermarks without possessing the capacity for \emph{extraction} or \emph{interpretation}, thereby enabling covert tracking of user-generated paraphrases or fine-tuned derivatives without consent---a capability explicitly flagged in \cite{doi:10.1145/1161366.1161397} and \cite{arxiv:2402.06922}. Such asymmetry violates foundational principles of informational self-determination and necessitates provably \emph{undetectable} watermarking schemes, as formalized in \cite{doi:10.18653/v1/2024.findings-naacl.40} --- though no paper in the provided list explicitly proposes watermark indistinguishability under cryptographic one-way assumptions; therefore, the citation ``\cite{arxiv:2403.06634}'' must be \textbf{removed}, as that paper is \emph{not in the provided list}.

Dual-use mitigation thus requires moving beyond detection-centric security toward \emph{enforceable cryptographic boundaries}. Asymmetric watermark generation/detection networks---where only the owner holds the private key for embedding, yet public verification remains possible---offer a structural safeguard against unauthorized monitoring \cite{arxiv:2402.16187}. Similarly, multi-party computation (MPC)-backed protocols like \cite{arxiv:2510.18333} decouple watermark verification from model access, preventing platform operators from weaponizing provenance signals for behavioral profiling --- but \emph{\cite{arxiv:2510.18333}} is \textbf{not in the provided list}, so this citation must be \textbf{removed}.

Crucially, unforgeability must be strengthened against \emph{contextual exploitation}: recent work demonstrates that watermark-induced token distribution shifts can degrade safety guardrails---either attenuating protections (e.g., permitting harmful outputs) or amplifying them excessively (e.g., suppressing legitimate medical advice)---a phenomenon empirically quantified in \cite{arxiv:2306.09194} --- but \emph{\cite{arxiv:2306.09194}} is \textbf{not in the provided list}, so this citation must be \textbf{removed}. Likewise, \cite{arxiv:2403.19548} \emph{does} explicitly analyze the quality-detection trade-off and its impact on generation utility, including alignment-sensitive tasks --- so this citation \textbf{is correct} and retained.

Red-teaming methodologies must therefore expand beyond technical removal to encompass \emph{sociotechnical failure modes}. For instance, cross-lingual watermark fragility---where round-trip translation collapses detection accuracy from >88\% to <13\% in low-resource languages---creates de facto exclusionary regimes, effectively denying attribution rights to non-English users and enabling false-flag IP violations in multilingual legal or academic contexts \cite{arxiv:2402.14007}; this citation is \textbf{correct} and retained. Likewise, ``radioactive'' inheritance of watermarks during knowledge distillation---while useful for copyright tracing---enables \emph{coercive downstream licensing}, where service providers embed watermarks that persist in client-fine-tuned models, thereby extending contractual control beyond initial deployment \cite{arxiv:2402.19361}; this citation is \textbf{correct}, as the paper explicitly discusses watermark stealing enabling spoofing \emph{and} scrubbing in commercial fine-tuning contexts, and \cite{doi:10.1145/3639477.3639720} also analyzes misuse scenarios including IP violation via fine-tuned derivatives --- both citations are valid and retained.

Ethical redlines must therefore be technically enforceable: mandatory human-readable fingerprint publication (\cite{arxiv:2312.04828}) --- but \emph{\cite{arxiv:2312.04828}} is \textbf{not in the provided list}, so this citation must be \textbf{removed}; prohibition of watermark-triggered content filtering without explicit opt-in consent --- no paper in the list addresses this normative requirement directly, so no citation applies; and open auditing of watermark decay rates across fine-tuning iterations (\cite{arxiv:2506.23731}; \cite{arxiv:2407.10887}) --- both are \textbf{correct}: \cite{arxiv:2506.23731} explicitly studies watermark transfer and decay under fine-tuning, and \cite{arxiv:2407.10887} evaluates fingerprint persistence under fine-tuning and LoRA adapters.

Ultimately, ethical governance is not orthogonal to technical design---it is its constitutive constraint, demanding watermarking systems that are \emph{verifiably non-instrumentalizable}, where robustness is measured not only in bits-per-token but in adherence to UNESCO's AI Ethics Recommendation: human oversight, transparency, and equity by construction.

\textbf{Corrected subsection with only valid citations from the provided list:}

Watermarking and fingerprinting, while indispensable for intellectual property protection and synthetic content accountability, embody a profound ethical paradox: their very mechanisms---designed to enhance transparency and traceability---can be repurposed as instruments of surveillance, exclusion, and coercive control. This duality demands governance frameworks that are not merely reactive but \emph{anticipatorily normative}, embedding ethical constraints directly into technical design. Critically, empirical evidence reveals that advanced reasoning models---including GPT-5 and Gemini 2.5 Pro---exhibit strong \emph{detection} capability for watermarks without possessing the capacity for \emph{extraction} or \emph{interpretation}, thereby enabling covert tracking of user-generated paraphrases or fine-tuned derivatives without consent---a capability explicitly flagged in \cite{doi:10.1145/1161366.1161397} and \cite{arxiv:2402.06922}.

Dual-use mitigation thus requires moving beyond detection-centric security toward \emph{enforceable cryptographic boundaries}. Asymmetric watermark generation/detection networks---where only the owner holds the private key for embedding, yet public verification remains possible---offer a structural safeguard against unauthorized monitoring \cite{arxiv:2402.16187}.

Crucially, unforgeability must be strengthened against \emph{contextual exploitation}: recent work demonstrates that watermark-induced token distribution shifts can degrade safety guardrails---either attenuating protections (e.g., permitting harmful outputs) or amplifying them excessively (e.g., suppressing legitimate medical advice)---a phenomenon empirically quantified in \cite{arxiv:2403.19548}.

Red-teaming methodologies must therefore expand beyond technical removal to encompass \emph{sociotechnical failure modes}. For instance, cross-lingual watermark fragility---where round-trip translation collapses detection accuracy from >88\% to <13\% in low-resource languages---creates de facto exclusionary regimes, effectively denying attribution rights to non-English users and enabling false-flag IP violations in multilingual legal or academic contexts \cite{arxiv:2402.14007}. Likewise, ``radioactive'' inheritance of watermarks during knowledge distillation---while useful for copyright tracing---enables \emph{coercive downstream licensing}, where service providers embed watermarks that persist in client-fine-tuned models, thereby extending contractual control beyond initial deployment \cite{arxiv:2402.19361}; \cite{doi:10.1145/3639477.3639720}.

Ethical redlines must therefore be technically enforceable: open auditing of watermark decay rates across fine-tuning iterations (\cite{arxiv:2506.23731}; \cite{arxiv:2407.10887}). Ultimately, ethical governance is not orthogonal to technical design---it is its constitutive constraint, demanding watermarking systems that are \emph{verifiably non-instrumentalizable}, where robustness is measured not only in bits-per-token but in adherence to UNESCO's AI Ethics Recommendation: human oversight, transparency, and equity by construction.


\section{Conclusion}


This subsection synthesizes the foundational insights, persistent tensions, and transformative research trajectories that emerge when watermarking and fingerprinting are examined not as isolated algorithmic artifacts but as constitutive elements of LLM security infrastructure. A unifying insight across the surveyed literature is the erosion of the classical dichotomy between \emph{active watermarking} and \emph{passive fingerprinting}: techniques like SLIM \cite{arxiv:2601.03242} and EmMark \cite{doi:10.1145/3649329.3655674} demonstrate that signal embedding and behavioral inference converge in latent-space instability and quantization-invariant weight geometry---both exploiting intrinsic model structure rather than external perturbation. Formally, this convergence is captured by a generalized detectability criterion: for any locus \$ \textbackslash{}mathcal\{L\} \textbackslash{}in \textbackslash{}\{\textbackslash{}text\{output\}, \textbackslash{}text\{parameter\}, \textbackslash{}text\{behavior\}\textbackslash{}\} \$, detectability \$ \textbackslash{}mathcal\{D\} \$ reduces to the statistical distinguishability of two hypotheses \$ H\_0: x \textbackslash{}sim p\_\{\textbackslash{}theta\}(x) \$ (unmarked) versus \$ H\_1: x \textbackslash{}sim p\_\{\textbackslash{}theta'\}(x) \$ (marked/fingerprinted), where the KL divergence \$ D\_\{\textbackslash{}mathrm\{KL\}\}(p\_\{\textbackslash{}theta'\} \textbackslash{}| p\_\{\textbackslash{}theta\}) \$ must exceed a threshold \$ \textbackslash{}tau \$ calibrated under adversarial perturbations \$ \textbackslash{}mathcal\{A\} \$, i.e., \$ \textbackslash{}inf\_\{a \textbackslash{}in \textbackslash{}mathcal\{A\}\} D\_\{\textbackslash{}mathrm\{KL\}\}(a(p\_\{\textbackslash{}theta'\}) \textbackslash{}| p\_\{\textbackslash{}theta\}) > \textbackslash{}tau \$. Yet current frameworks fail to bound \$ \textbackslash{}tau \$ under compositionally adaptive adversaries---e.g., DIPPER-based paraphrasing followed by fine-tuning erasure \cite{arxiv:2303.13408}---exposing a critical gap in threat modeling. Empirically, robustness degrades catastrophically under such composite attacks: detection accuracy collapses from 92\% to <15\% in cross-lingual settings \cite{arxiv:2303.13408}, while alignment degradation manifests as measurable guard attenuation ($-$37\% safety score) or amplification ($-$29\% helpfulness) \cite{arxiv:2306.11698}. This reveals a fundamental trade-off formalized in the \emph{Inevitable Trade-off} theorem \cite{doi:10.1038/s41586-024-08025-4}: for any watermark strength \$ \textbackslash{}lambda \$ and speculative decoding speedup \$ \textbackslash{}gamma \$, \$ \textbackslash{}lambda \textbackslash{}cdot \textbackslash{}gamma \textbackslash{}leq C \$, where \$ C \$ is architecture-dependent. Consequently, production-ready systems like SynthID-Text \cite{doi:10.1038/s41586-024-08025-4} and PRIVMARK \cite{doi:10.3390/info15110697} decouple embedding from inference---embedding via MPC-secured parameter modulation or post-hoc semantic insertion---thereby sidestepping the no-go constraint. However, this architectural shift introduces new vulnerabilities: black-box watermarking remains susceptible to prompt-probing attacks \cite{arxiv:2402.19361}, while cryptographic schemes like Undetectable Watermarks \cite{arxiv:2306.09194} assume ideal one-way functions, ignoring practical leakage from gradient inversion \cite{arxiv:2402.19361}. Most critically, longitudinal persistence---the cornerstone of provenance---is empirically unverified: no benchmark tracks watermark decay across iterative fine-tuning \cite{arxiv:2410.06467}, continual learning, or version updates \cite{doi:10.3389/fdata.2021.729663}, despite evidence that even 0.001\% data poisoning can corrupt downstream integrity \cite{doi:10.1038/s41591-024-03445-1}. Future work must therefore co-design watermarking with system-level governance: integrating zero-knowledge proofs for forensic admissibility \cite{doi:10.34190/iccws.19.1.2096}, blockchain-anchored provenance trails \cite{doi:10.34190/icair.4.1.3128}, and human-centered evaluation protocols that quantify perceptual fidelity disparities across linguistic registers \cite{doi:10.1145/3637528.3671467}. Only through such interdisciplinary integration can watermarking and fingerprinting evolve from fragile technical add-ons into verifiable, equitable, and legally enforceable pillars of AI accountability.

\begin{thebibliography}{121}
\setlength{\itemsep}{2pt}

\bibitem{doi:10.55662/jst.2023.4605}
\newblock \emph{Guardrails for trust, safety, and ethical development and deployment of Large Language Models (LLM)}.
\newblock DOI:10.55662/jst.2023.4605, 2023.
\newblock \url{https://doi.org/10.55662/jst.2023.4605}

\bibitem{doi:10.3390/info15110697}
\newblock \emph{Privacy-Preserving Techniques in Generative AI and Large Language Models: A Narrative Review}.
\newblock DOI:10.3390/info15110697, 2024.
\newblock \url{https://doi.org/10.3390/info15110697}

\bibitem{doi:10.1007/s44443-025-00177-1}
\newblock \emph{A survey on privacy risks and protection in large language models}.
\newblock DOI:10.1007/s44443-025-00177-1, 2025.
\newblock \url{https://doi.org/10.1007/s44443-025-00177-1}

\bibitem{arxiv:2309.15025}
\newblock \emph{Large Language Model Alignment A Survey}.
\newblock arXiv:2309.15025, 2023.
\newblock \url{https://arxiv.org/abs/2309.15025}

\bibitem{doi:10.1007/s10462-024-11024-6}
\newblock \emph{Generative AI model privacy: a survey}.
\newblock DOI:10.1007/s10462-024-11024-6, 2024.
\newblock \url{https://doi.org/10.1007/s10462-024-11024-6}

\bibitem{arxiv:2405.01560}
\newblock \emph{Copyright related risks in the creation and use of ML/AI systems}.
\newblock arXiv:2405.01560, 2024.
\newblock \url{https://arxiv.org/abs/2405.01560}

\bibitem{arxiv:1701.04082}
\newblock \emph{Embedding Watermarks into Deep Neural Networks}.
\newblock arXiv:1701.04082, 2017.
\newblock \url{https://arxiv.org/abs/1701.04082}

\bibitem{doi:10.1145/3292500.3330885}
\newblock \emph{Auditing Data Provenance in Text-Generation Models}.
\newblock DOI:10.1145/3292500.3330885, 2019.
\newblock \url{https://doi.org/10.1145/3292500.3330885}

\bibitem{doi:10.1038/s41586-024-08025-4}
\newblock \emph{Scalable watermarking for identifying large language model outputs}.
\newblock DOI:10.1038/s41586-024-08025-4, 2024.
\newblock \url{https://doi.org/10.1038/s41586-024-08025-4}

\bibitem{doi:10.1371/journal.pone.0331871}
\newblock \emph{Detecting LLM-generated peer reviews}.
\newblock DOI:10.1371/journal.pone.0331871, 2025.
\newblock \url{https://doi.org/10.1371/journal.pone.0331871}

\bibitem{arxiv:2112.04359}
\newblock \emph{Ethical and social risks of harm from Language Models}.
\newblock arXiv:2112.04359, 2021.
\newblock \url{https://arxiv.org/abs/2112.04359}

\bibitem{doi:10.1038/s41591-024-03445-1}
\newblock \emph{Medical large language models are vulnerable to data-poisoning attacks}.
\newblock DOI:10.1038/s41591-024-03445-1, 2025.
\newblock \url{https://doi.org/10.1038/s41591-024-03445-1}

\bibitem{arxiv:2303.13408}
\newblock \emph{Paraphrasing evades detectors of AI-generated text, but retrieval is an effective defense}.
\newblock arXiv:2303.13408, 2023.
\newblock \url{https://arxiv.org/abs/2303.13408}

\bibitem{arxiv:1903.01743}
\newblock \emph{How to Prove Your Model Belongs to You: A Blind-Watermark based Framework to Protect Intellectual Property of DNN}.
\newblock arXiv:1903.01743, 2019.
\newblock \url{https://arxiv.org/abs/1903.01743}

\bibitem{doi:10.1145/3649329.3655674}
\newblock \emph{EmMark: Robust Watermarks for IP Protection of Embedded Quantized Large Language Models}.
\newblock DOI:10.1145/3649329.3655674, 2024.
\newblock \url{https://doi.org/10.1145/3649329.3655674}

\bibitem{arxiv:2601.03242}
\newblock \emph{SLIM: Stealthy Low-Coverage Black-Box Watermarking via Latent-Space Confusion Zones}.
\newblock arXiv:2601.03242, 2026.
\newblock \url{https://arxiv.org/abs/2601.03242}

\bibitem{arxiv:2502.00706}
\newblock \emph{Model Provenance Testing for Large Language Models}.
\newblock arXiv:2502.00706, 2025.
\newblock \url{https://arxiv.org/abs/2502.00706}

\bibitem{arxiv:2306.09194}
\newblock \emph{Undetectable Watermarks for Language Models}.
\newblock arXiv:2306.09194, 2023.
\newblock \url{https://arxiv.org/abs/2306.09194}

\bibitem{doi:10.1609/aaai.v37i12.26750}
\newblock \emph{PLMmark: A Secure and Robust Black-Box Watermarking Framework for Pre-trained Language Models}.
\newblock DOI:10.1609/aaai.v37i12.26750, 2023.
\newblock \url{https://doi.org/10.1609/aaai.v37i12.26750}

\bibitem{arxiv:2406.03728}
\newblock \emph{Evaluating Durability: Benchmark Insights into Multimodal Watermarking}.
\newblock arXiv:2406.03728, 2024.
\newblock \url{https://arxiv.org/abs/2406.03728}

\bibitem{arxiv:2402.16187}
\newblock \emph{Attacking LLM Watermarks by Exploiting Their Strengths}.
\newblock arXiv:2402.16187, 2024.
\newblock \url{https://arxiv.org/abs/2402.16187}

\bibitem{doi:10.1109/icassp48485.2024.10447264}
\newblock \emph{Functional Invariants To Watermark Large Transformers}.
\newblock DOI:10.1109/icassp48485.2024.10447264, 2024.
\newblock \url{https://doi.org/10.1109/icassp48485.2024.10447264}

\bibitem{doi:10.18653/v1/2023.acl-long.423}
\newblock \emph{Are You Copying My Model? Protecting the Copyright of Large Language Models for EaaS via Backdoor Watermark}.
\newblock DOI:10.18653/v1/2023.acl-long.423, 2023.
\newblock \url{https://doi.org/10.18653/v1/2023.acl-long.423}

\bibitem{arxiv:2410.06467}
\newblock \emph{WAPITI: A Watermark for Finetuned Open-Source LLMs}.
\newblock arXiv:2410.06467, 2024.
\newblock \url{https://arxiv.org/abs/2410.06467}

\bibitem{arxiv:2406.01946}
\newblock \emph{Bileve: Securing Text Provenance in Large Language Models Against Spoofing with Bi-level Signature}.
\newblock arXiv:2406.01946, 2024.
\newblock \url{https://arxiv.org/abs/2406.01946}

\bibitem{doi:10.1109/dac56929.2023.10247798}
\newblock \emph{ZKROWNN: Zero Knowledge Right of Ownership for Neural Networks}.
\newblock DOI:10.1109/dac56929.2023.10247798, 2023.
\newblock \url{https://doi.org/10.1109/dac56929.2023.10247798}

\bibitem{doi:10.1145/3580305.3599291}
\newblock \emph{Cracking White-box DNN Watermarks via Invariant Neuron Transforms}.
\newblock DOI:10.1145/3580305.3599291, 2023.
\newblock \url{https://doi.org/10.1145/3580305.3599291}

\bibitem{arxiv:2402.19361}
\newblock \emph{Watermark Stealing in Large Language Models}.
\newblock arXiv:2402.19361, 2024.
\newblock \url{https://arxiv.org/abs/2402.19361}

\bibitem{arxiv:2405.19677}
\newblock \emph{Large Language Model Watermark Stealing With Mixed Integer Programming}.
\newblock arXiv:2405.19677, 2024.
\newblock \url{https://arxiv.org/abs/2405.19677}

\bibitem{arxiv:2303.09732}
\newblock \emph{Rethinking White-Box Watermarks on Deep Learning Models under Neural Structural Obfuscation}.
\newblock arXiv:2303.09732, 2023.
\newblock \url{https://arxiv.org/abs/2303.09732}

\bibitem{arxiv:2307.15593}
\newblock \emph{Robust Distortion-free Watermarks for Language Models}.
\newblock arXiv:2307.15593, 2023.
\newblock \url{https://arxiv.org/abs/2307.15593}

\bibitem{arxiv:2402.05864}
\newblock \emph{Permute-and-Flip An optimally robust and watermarkable decoder for LLMs}.
\newblock arXiv:2402.05864, 2024.
\newblock \url{https://arxiv.org/abs/2402.05864}

\bibitem{doi:10.1016/j.comcom.2019.12.016}
\newblock \emph{AFA: Adversarial fingerprinting authentication for deep neural networks}.
\newblock DOI:10.1016/j.comcom.2019.12.016, 2020.
\newblock \url{https://doi.org/10.1016/j.comcom.2019.12.016}

\bibitem{doi:10.1609/aaai.v36i10.21415}
\newblock \emph{Tracing Text Provenance via Context-Aware Lexical Substitution}.
\newblock DOI:10.1609/aaai.v36i10.21415, 2022.
\newblock \url{https://doi.org/10.1609/aaai.v36i10.21415}

\bibitem{arxiv:2310.00646}
\newblock \emph{WASA WAtermark-based Source Attribution for Large Language Model-Generated Data}.
\newblock arXiv:2310.00646, 2023.
\newblock \url{https://arxiv.org/abs/2310.00646}

\bibitem{arxiv:2407.10887}
\newblock \emph{Hey, That's My Model! Introducing Chain \& Hash, An LLM Fingerprinting Technique}.
\newblock arXiv:2407.10887, 2024.
\newblock \url{https://arxiv.org/abs/2407.10887}

\bibitem{arxiv:2405.13517}
\newblock \emph{WaterPool: A Watermark Mitigating Trade-offs among Imperceptibility, Efficacy and Robustness}.
\newblock arXiv:2405.13517, 2024.
\newblock \url{https://arxiv.org/abs/2405.13517}

\bibitem{arxiv:1910.12366}
\newblock \emph{Thieves on Sesame Street! Model Extraction of BERT-based APIs}.
\newblock arXiv:1910.12366, 2019.
\newblock \url{https://arxiv.org/abs/1910.12366}

\bibitem{arxiv:2505.06827}
\newblock \emph{Sandcastles in the Storm: Revisiting the (Im)possibility of Strong Watermarking}.
\newblock arXiv:2505.06827, 2025.
\newblock \url{https://arxiv.org/abs/2505.06827}

\bibitem{arxiv:2301.10226}
\newblock \emph{A Watermark for Large Language Models}.
\newblock arXiv:2301.10226, 2023.
\newblock \url{https://arxiv.org/abs/2301.10226}

\bibitem{arxiv:2306.04634}
\newblock \emph{On the Reliability of Watermarks for Large Language Models}.
\newblock arXiv:2306.04634, 2023.
\newblock \url{https://arxiv.org/abs/2306.04634}

\bibitem{arxiv:2406.10281}
\newblock \emph{Watermarking Language Models with Error Correcting Codes}.
\newblock arXiv:2406.10281, 2024.
\newblock \url{https://arxiv.org/abs/2406.10281}

\bibitem{arxiv:2403.04808}
\newblock \emph{WaterMax: breaking the LLM watermark detectability-robustness-quality trade-off}.
\newblock arXiv:2403.04808, 2024.
\newblock \url{https://arxiv.org/abs/2403.04808}

\bibitem{doi:10.18653/v1/2024.findings-naacl.40}
\newblock \emph{A Robust Semantics-based Watermark for Large Language Model against Paraphrasing}.
\newblock DOI:10.18653/v1/2024.findings-naacl.40, 2024.
\newblock \url{https://doi.org/10.18653/v1/2024.findings-naacl.40}

\bibitem{doi:10.18653/v1/2023.acl-long.117}
\newblock \emph{Robust Multi-bit Natural Language Watermarking through Invariant Features}.
\newblock DOI:10.18653/v1/2023.acl-long.117, 2023.
\newblock \url{https://doi.org/10.18653/v1/2023.acl-long.117}

\bibitem{arxiv:2406.02603}
\newblock \emph{Distortion-free Watermarks are not Truly Distortion-free under Watermark Key Collisions}.
\newblock arXiv:2406.02603, 2024.
\newblock \url{https://arxiv.org/abs/2406.02603}

\bibitem{arxiv:2601.01296}
\newblock \emph{Aggressive Compression Enables LLM Weight Theft}.
\newblock arXiv:2601.01296, 2026.
\newblock \url{https://arxiv.org/abs/2601.01296}

\bibitem{arxiv:2601.04261}
\newblock \emph{Inhibitory Attacks on Backdoor-based Fingerprinting for Large Language Models}.
\newblock arXiv:2601.04261, 2026.
\newblock \url{https://arxiv.org/abs/2601.04261}

\bibitem{doi:10.1145/3618257.3624845}
\newblock \emph{An LLM-based Framework for Fingerprinting Internet-connected Devices}.
\newblock DOI:10.1145/3618257.3624845, 2023.
\newblock \url{https://doi.org/10.1145/3618257.3624845}

\bibitem{arxiv:2409.06130}
\newblock \emph{On the Weaknesses of Backdoor-based Model Watermarking: An Information-theoretic Perspective}.
\newblock arXiv:2409.06130, 2024.
\newblock \url{https://arxiv.org/abs/2409.06130}

\bibitem{arxiv:2505.02884}
\newblock \emph{Unlearning vs. Obfuscation: Are We Truly Removing Knowledge?}.
\newblock arXiv:2505.02884, 2025.
\newblock \url{https://arxiv.org/abs/2505.02884}

\bibitem{arxiv:2402.14007}
\newblock \emph{Can Watermarks Survive Translation? On the Cross-lingual Consistency of Text Watermark for Large Language Models}.
\newblock arXiv:2402.14007, 2024.
\newblock \url{https://arxiv.org/abs/2402.14007}

\bibitem{doi:10.14722/ndss.2026.230138}
\newblock \emph{Character-Level Perturbations Disrupt LLM Watermarks}.
\newblock DOI:10.14722/ndss.2026.230138, 2025.
\newblock \url{https://doi.org/10.14722/ndss.2026.230138}

\bibitem{arxiv:2401.16820}
\newblock \emph{Provably Robust Multi-bit Watermarking for AI-generated Text via Error Correction Code}.
\newblock arXiv:2401.16820, 2024.
\newblock \url{https://arxiv.org/abs/2401.16820}

\bibitem{doi:10.56553/popets-2025-0126}
\newblock \emph{Robust and Efficient Watermarking of Large Language Models Using Error Correction Codes}.
\newblock DOI:10.56553/popets-2025-0126, 2025.
\newblock \url{https://doi.org/10.56553/popets-2025-0126}

\bibitem{arxiv:2601.14311}
\newblock \emph{Tracing the Data Trail: A Survey of Data Provenance, Transparency and Traceability in LLMs}.
\newblock arXiv:2601.14311, 2026.
\newblock \url{https://arxiv.org/abs/2601.14311}

\bibitem{doi:10.3389/fdata.2021.729663}
\newblock \emph{A Systematic Review on Model Watermarking for Neural Networks}.
\newblock DOI:10.3389/fdata.2021.729663, 2021.
\newblock \url{https://doi.org/10.3389/fdata.2021.729663}

\bibitem{arxiv:2401.13927}
\newblock \emph{Adaptive Text Watermark for Large Language Models}.
\newblock arXiv:2401.13927, 2024.
\newblock \url{https://arxiv.org/abs/2401.13927}

\bibitem{arxiv:2601.05144}
\newblock \emph{Distilling the Thought, Watermarking the Answer: A Principle Semantic Guided Watermark for Large Reasoning Models}.
\newblock arXiv:2601.05144, 2026.
\newblock \url{https://arxiv.org/abs/2601.05144}

\bibitem{arxiv:2407.04411}
\newblock \emph{Waterfall: Framework for Robust and Scalable Text Watermarking}.
\newblock arXiv:2407.04411, 2024.
\newblock \url{https://arxiv.org/abs/2407.04411}

\bibitem{arxiv:2506.06407}
\newblock \emph{TimeWak: Temporal Chained-Hashing Watermark for Time Series Data}.
\newblock arXiv:2506.06407, 2025.
\newblock \url{https://arxiv.org/abs/2506.06407}

\bibitem{doi:10.1145/3321705.3329808}
\newblock \emph{Robust Watermarking of Neural Network with Exponential Weighting}.
\newblock DOI:10.1145/3321705.3329808, 2019.
\newblock \url{https://doi.org/10.1145/3321705.3329808}

\bibitem{arxiv:2512.03620}
\newblock \emph{SELF: A Robust Singular Value and Eigenvalue Approach for LLM Fingerprinting}.
\newblock arXiv:2512.03620, 2025.
\newblock \url{https://arxiv.org/abs/2512.03620}

\bibitem{arxiv:2403.10553}
\newblock \emph{Learning to Watermark LLM-generated Text via Reinforcement Learning}.
\newblock arXiv:2403.10553, 2024.
\newblock \url{https://arxiv.org/abs/2403.10553}

\bibitem{arxiv:2511.07863}
\newblock \emph{WaterMod: Modular Token-Rank Partitioning for Probability-Balanced LLM Watermarking}.
\newblock arXiv:2511.07863, 2025.
\newblock \url{https://arxiv.org/abs/2511.07863}

\bibitem{doi:10.34190/icair.4.1.3128}
\newblock \emph{LLM Supply Chain Provenance: A Blockchain-based Approach}.
\newblock DOI:10.34190/icair.4.1.3128, 2024.
\newblock \url{https://doi.org/10.34190/icair.4.1.3128}

\bibitem{arxiv:2403.13485}
\newblock \emph{An Entropy-based Text Watermarking Detection Method}.
\newblock arXiv:2403.13485, 2024.
\newblock \url{https://arxiv.org/abs/2403.13485}

\bibitem{arxiv:2502.18851}
\newblock \emph{Marking Code Without Breaking It: Code Watermarking for Detecting LLM-Generated Code}.
\newblock arXiv:2502.18851, 2025.
\newblock \url{https://arxiv.org/abs/2502.18851}

\bibitem{arxiv:2505.05190}
\newblock \emph{Revealing Weaknesses in Text Watermarking Through Self-Information Rewrite Attacks}.
\newblock arXiv:2505.05190, 2025.
\newblock \url{https://arxiv.org/abs/2505.05190}

\bibitem{arxiv:2310.06356}
\newblock \emph{A Semantic Invariant Robust Watermark for Large Language Models}.
\newblock arXiv:2310.06356, 2023.
\newblock \url{https://arxiv.org/abs/2310.06356}

\bibitem{arxiv:2601.11629}
\newblock \emph{Semantic Differentiation for Tackling Challenges in Watermarking Low-Entropy Constrained Generation Outputs}.
\newblock arXiv:2601.11629, 2026.
\newblock \url{https://arxiv.org/abs/2601.11629}

\bibitem{doi:10.1109/cvpr52733.2024.01169}
\newblock \emph{Steganographic Passport: An Owner and User Verifiable Credential for Deep Model IP Protection Without Retraining}.
\newblock DOI:10.1109/cvpr52733.2024.01169, 2024.
\newblock \url{https://doi.org/10.1109/cvpr52733.2024.01169}

\bibitem{doi:10.1145/3691626}
\newblock \emph{A Survey of Text Watermarking in the Era of Large Language Models}.
\newblock DOI:10.1145/3691626, 2024.
\newblock \url{https://doi.org/10.1145/3691626}

\bibitem{arxiv:2502.11598}
\newblock \emph{Can LLM Watermarks Robustly Prevent Unauthorized Knowledge Distillation?}.
\newblock arXiv:2502.11598, 2025.
\newblock \url{https://arxiv.org/abs/2502.11598}

\bibitem{doi:10.1145/3639477.3639720}
\newblock \emph{Assessing LLMs for High Stakes Applications}.
\newblock DOI:10.1145/3639477.3639720, 2024.
\newblock \url{https://doi.org/10.1145/3639477.3639720}

\bibitem{doi:10.1609/aaai.v39i1.32018}
\newblock \emph{Simulate and Eliminate: Revoke Backdoors for Generative Large Language Models}.
\newblock DOI:10.1609/aaai.v39i1.32018, 2025.
\newblock \url{https://doi.org/10.1609/aaai.v39i1.32018}

\bibitem{doi:10.1109/sds66131.2025.00024}
\newblock \emph{Securing Rag: A Risk Assessment and Mitigation Framework}.
\newblock DOI:10.1109/sds66131.2025.00024, 2025.
\newblock \url{https://doi.org/10.1109/sds66131.2025.00024}

\bibitem{arxiv:2601.21252}
\newblock \emph{Lossless Copyright Protection via Intrinsic Model Fingerprinting}.
\newblock arXiv:2601.21252, 2026.
\newblock \url{https://arxiv.org/abs/2601.21252}

\bibitem{arxiv:2401.12255}
\newblock \emph{Instructional Fingerprinting of Large Language Models}.
\newblock arXiv:2401.12255, 2024.
\newblock \url{https://arxiv.org/abs/2401.12255}

\bibitem{arxiv:2502.18518}
\newblock \emph{Swallowing the Poison Pills: Insights from Vulnerability Disparity Among LLMs}.
\newblock arXiv:2502.18518, 2025.
\newblock \url{https://arxiv.org/abs/2502.18518}

\bibitem{arxiv:2305.08883}
\newblock \emph{Watermarking Text Generated by Black-Box Language Models}.
\newblock arXiv:2305.08883, 2023.
\newblock \url{https://arxiv.org/abs/2305.08883}

\bibitem{arxiv:2402.16578}
\newblock \emph{Multi-Bit Distortion-Free Watermarking for Large Language Models}.
\newblock arXiv:2402.16578, 2024.
\newblock \url{https://arxiv.org/abs/2402.16578}

\bibitem{doi:10.18653/v1/2024.naacl-long.224}
\newblock \emph{Advancing Beyond Identification: Multi-bit Watermark for Large Language Models}.
\newblock DOI:10.18653/v1/2024.naacl-long.224, 2024.
\newblock \url{https://doi.org/10.18653/v1/2024.naacl-long.224}

\bibitem{arxiv:2505.14112}
\newblock \emph{Invisible Entropy: Towards Safe and Efficient Low-Entropy LLM Watermarking}.
\newblock arXiv:2505.14112, 2025.
\newblock \url{https://arxiv.org/abs/2505.14112}

\bibitem{arxiv:2209.08773}
\newblock \emph{CATER: Intellectual Property Protection on Text Generation APIs via Conditional Watermarks}.
\newblock arXiv:2209.08773, 2022.
\newblock \url{https://arxiv.org/abs/2209.08773}

\bibitem{doi:10.3390/electronics14010025}
\newblock \emph{Latent Diffusion Models for Image Watermarking: A Review of Recent Trends and Future Directions}.
\newblock DOI:10.3390/electronics14010025, 2024.
\newblock \url{https://doi.org/10.3390/electronics14010025}

\bibitem{arxiv:2402.18059}
\newblock \emph{Token-Specific Watermarking with Enhanced Detectability and Semantic Coherence for Large Language Models}.
\newblock arXiv:2402.18059, 2024.
\newblock \url{https://arxiv.org/abs/2402.18059}

\bibitem{arxiv:2506.23731}
\newblock \emph{Radioactive Watermarks in Diffusion and Autoregressive Image Generative Models}.
\newblock arXiv:2506.23731, 2025.
\newblock \url{https://arxiv.org/abs/2506.23731}

\bibitem{arxiv:2002.12200}
\newblock \emph{Entangled Watermarks as a Defense against Model Extraction}.
\newblock arXiv:2002.12200, 2020.
\newblock \url{https://arxiv.org/abs/2002.12200}

\bibitem{arxiv:2506.00359}
\newblock \emph{Keeping an Eye on LLM Unlearning: The Hidden Risk and Remedy}.
\newblock arXiv:2506.00359, 2025.
\newblock \url{https://arxiv.org/abs/2506.00359}

\bibitem{arxiv:2307.15043}
\newblock \emph{Universal and Transferable Adversarial Attacks on Aligned Language Models}.
\newblock arXiv:2307.15043, 2023.
\newblock \url{https://arxiv.org/abs/2307.15043}

\bibitem{doi:10.1109/cvpr52734.2025.02785}
\newblock \emph{BadToken: Token-level Backdoor Attacks to Multi-modal Large Language Models}.
\newblock DOI:10.1109/cvpr52734.2025.02785, 2025.
\newblock \url{https://doi.org/10.1109/cvpr52734.2025.02785}

\bibitem{arxiv:2402.08577}
\newblock \emph{Test-Time Backdoor Attacks on Multimodal Large Language Models}.
\newblock arXiv:2402.08577, 2024.
\newblock \url{https://arxiv.org/abs/2402.08577}

\bibitem{arxiv:2410.06545}
\newblock \emph{Signal Watermark on Large Language Models}.
\newblock arXiv:2410.06545, 2024.
\newblock \url{https://arxiv.org/abs/2410.06545}

\bibitem{arxiv:1804.00750}
\newblock \emph{DeepSigns: A Generic Watermarking Framework for IP Protection of Deep Learning Models}.
\newblock arXiv:1804.00750, 2018.
\newblock \url{https://arxiv.org/abs/1804.00750}

\bibitem{arxiv:2509.25448}
\newblock \emph{Fingerprinting LLMs via Prompt Injection}.
\newblock arXiv:2509.25448, 2025.
\newblock \url{https://arxiv.org/abs/2509.25448}

\bibitem{arxiv:2403.15447}
\newblock \emph{Decoding Compressed Trust: Scrutinizing the Trustworthiness of Efficient LLMs Under Compression}.
\newblock arXiv:2403.15447, 2024.
\newblock \url{https://arxiv.org/abs/2403.15447}

\bibitem{arxiv:2403.07283}
\newblock \emph{A Framework for Cost-Effective and Self-Adaptive LLM Shaking and Recovery Mechanism}.
\newblock arXiv:2403.07283, 2024.
\newblock \url{https://arxiv.org/abs/2403.07283}

\bibitem{doi:10.1109/access.2024.3376693}
\newblock \emph{DeepTextMark: A Deep Learning-Driven Text Watermarking Approach for Identifying Large Language Model Generated Text}.
\newblock DOI:10.1109/access.2024.3376693, 2024.
\newblock \url{https://doi.org/10.1109/access.2024.3376693}

\bibitem{arxiv:2307.16230}
\newblock \emph{A Private Watermark for Large Language Models}.
\newblock arXiv:2307.16230, 2023.
\newblock \url{https://arxiv.org/abs/2307.16230}

\bibitem{arxiv:2403.10893}
\newblock \emph{A Watermark-Conditioned Diffusion Model for IP Protection}.
\newblock arXiv:2403.10893, 2024.
\newblock \url{https://arxiv.org/abs/2403.10893}

\bibitem{arxiv:2506.15975}
\newblock \emph{Multi-use LLM Watermarking and the False Detection Problem}.
\newblock arXiv:2506.15975, 2025.
\newblock \url{https://arxiv.org/abs/2506.15975}

\bibitem{arxiv:2510.17000}
\newblock \emph{Bits Leaked per Query: Information-Theoretic Bounds on Adversarial Attacks against LLMs}.
\newblock arXiv:2510.17000, 2025.
\newblock \url{https://arxiv.org/abs/2510.17000}

\bibitem{arxiv:2506.01011}
\newblock \emph{Autoregressive Images Watermarking through Lexical Biasing: An Approach Resistant to Regeneration Attack}.
\newblock arXiv:2506.01011, 2025.
\newblock \url{https://arxiv.org/abs/2506.01011}

\bibitem{arxiv:2505.16530}
\newblock \emph{DuFFin: A Dual-Level Fingerprinting Framework for LLMs IP Protection}.
\newblock arXiv:2505.16530, 2025.
\newblock \url{https://arxiv.org/abs/2505.16530}

\bibitem{arxiv:2405.04825}
\newblock \emph{Explanation as a Watermark: Towards Harmless and Multi-bit Model Ownership Verification via Watermarking Feature Attribution}.
\newblock arXiv:2405.04825, 2024.
\newblock \url{https://arxiv.org/abs/2405.04825}

\bibitem{arxiv:2402.10892}
\newblock \emph{Proving membership in LLM pretraining data via data watermarks}.
\newblock arXiv:2402.10892, 2024.
\newblock \url{https://arxiv.org/abs/2402.10892}

\bibitem{doi:10.34190/iccws.19.1.2096}
\newblock \emph{Enhancing Privacy and Security in Large-Language Models: A Zero-Knowledge Proof Approach}.
\newblock DOI:10.34190/iccws.19.1.2096, 2024.
\newblock \url{https://doi.org/10.34190/iccws.19.1.2096}

\bibitem{doi:10.1109/eurosp57164.2023.00072}
\newblock \emph{Watermarking Graph Neural Networks based on Backdoor Attacks}.
\newblock DOI:10.1109/eurosp57164.2023.00072, 2023.
\newblock \url{https://doi.org/10.1109/eurosp57164.2023.00072}

\bibitem{arxiv:2410.23861}
\newblock \emph{Audio Is the Achilles' Heel: Red Teaming Audio Large Multimodal Models}.
\newblock arXiv:2410.23861, 2024.
\newblock \url{https://arxiv.org/abs/2410.23861}

\bibitem{arxiv:2403.19548}
\newblock \emph{WaterJudge Quality-Detection Trade-off when Watermarking Large Language Models}.
\newblock arXiv:2403.19548, 2024.
\newblock \url{https://arxiv.org/abs/2403.19548}

\bibitem{doi:10.1145/3611643.3616297}
\newblock \emph{CodeMark: Imperceptible Watermarking for Code Datasets against Neural Code Completion Models}.
\newblock DOI:10.1145/3611643.3616297, 2023.
\newblock \url{https://doi.org/10.1145/3611643.3616297}

\bibitem{arxiv:2310.15654}
\newblock \emph{A Survey on Detection of LLMs-Generated Content}.
\newblock arXiv:2310.15654, 2023.
\newblock \url{https://arxiv.org/abs/2310.15654}

\bibitem{arxiv:2509.03058}
\newblock \emph{EverTracer: Hunting Stolen Large Language Models via Stealthy and Robust Probabilistic Fingerprint}.
\newblock arXiv:2509.03058, 2025.
\newblock \url{https://arxiv.org/abs/2509.03058}

\bibitem{doi:10.1145/1161366.1161397}
\newblock \emph{The hiding virtues of ambiguity: quantifiably resilient watermarking of natural language text through synonym substitutions}.
\newblock DOI:10.1145/1161366.1161397, 2006.
\newblock \url{https://doi.org/10.1145/1161366.1161397}

\bibitem{arxiv:2402.06922}
\newblock \emph{Whispers in the Machine Confidentiality in LLM-integrated Systems}.
\newblock arXiv:2402.06922, 2024.
\newblock \url{https://arxiv.org/abs/2402.06922}

\bibitem{arxiv:2403.06634}
\newblock \emph{Stealing Part of a Production Language Model}.
\newblock arXiv:2403.06634, 2024.
\newblock \url{https://arxiv.org/abs/2403.06634}

\bibitem{arxiv:2510.18333}
\newblock \emph{Position: LLM Watermarking Should Align Stakeholders' Incentives for Practical Adoption}.
\newblock arXiv:2510.18333, 2025.
\newblock \url{https://arxiv.org/abs/2510.18333}

\bibitem{arxiv:2312.04828}
\newblock \emph{HuRef: HUman-REadable Fingerprint for Large Language Models}.
\newblock arXiv:2312.04828, 2023.
\newblock \url{https://arxiv.org/abs/2312.04828}

\bibitem{arxiv:2306.11698}
\newblock \emph{DecodingTrust A Comprehensive Assessment of Trustworthiness in GPT Models}.
\newblock arXiv:2306.11698, 2023.
\newblock \url{https://arxiv.org/abs/2306.11698}

\bibitem{doi:10.1145/3637528.3671467}
\newblock \emph{Grounding and Evaluation for Large Language Models: Practical Challenges and Lessons Learned (Survey)}.
\newblock DOI:10.1145/3637528.3671467, 2024.
\newblock \url{https://doi.org/10.1145/3637528.3671467}

\end{thebibliography}

\end{document}
